\documentclass[aps,preprint]{revtex4}%
\usepackage{amsfonts}
\usepackage{amsmath}
\usepackage{amssymb}
\usepackage{graphicx}%
\setcounter{MaxMatrixCols}{30}
%TCIDATA{OutputFilter=latex2.dll}
%TCIDATA{Version=5.50.0.2953}
%TCIDATA{CSTFile=revtex4.cst}
%TCIDATA{Created=Monday, September 20, 2004 09:58:21}
%TCIDATA{LastRevised=Thursday, September 11, 2008 13:25:55}
%TCIDATA{<META NAME="GraphicsSave" CONTENT="32">}
%TCIDATA{<META NAME="SaveForMode" CONTENT="1">}
%TCIDATA{BibliographyScheme=BibTeX}
%TCIDATA{<META NAME="DocumentShell" CONTENT="Articles\SW\REVTeX 4">}
%TCIDATA{Language=American English}
%BeginMSIPreambleData
\providecommand{\U}[1]{\protect\rule{.1in}{.1in}}
%EndMSIPreambleData
\newtheorem{theorem}{Theorem}
\newtheorem{acknowledgement}[theorem]{Acknowledgement}
\newtheorem{algorithm}[theorem]{Algorithm}
\newtheorem{axiom}[theorem]{Axiom}
\newtheorem{claim}[theorem]{Claim}
\newtheorem{conclusion}[theorem]{Conclusion}
\newtheorem{condition}[theorem]{Condition}
\newtheorem{conjecture}[theorem]{Conjecture}
\newtheorem{corollary}[theorem]{Corollary}
\newtheorem{criterion}[theorem]{Criterion}
\newtheorem{definition}[theorem]{Definition}
\newtheorem{example}[theorem]{Example}
\newtheorem{exercise}[theorem]{Exercise}
\newtheorem{lemma}[theorem]{Lemma}
\newtheorem{notation}[theorem]{Notation}
\newtheorem{problem}[theorem]{Problem}
\newtheorem{proposition}[theorem]{Proposition}
\newtheorem{remark}[theorem]{Remark}
\newtheorem{solution}[theorem]{Solution}
\newtheorem{summary}[theorem]{Summary}
\newenvironment{proof}[1][Proof]{\noindent\textbf{#1.} }{\ \rule{0.5em}{0.5em}}
\begin{document}
\preprint{ }
\title[Janaks Theorem]{Janaks Theorem and Orbital Energy Functionals (in progress)}
\author{B. Weiner}
\affiliation{Department of Physics, Pennsylvania State University}
\keywords{}
\pacs{}

\begin{abstract}
Janaks theorem shows that the rate of change of the total N-electron energy
with respect to the occupation numbers of the natural orbitals of the model
Kohn-Sham (KS) independent particle state are equal to the Lagrange
multipliers associated with the normalization constraint imposed on these
orbitals. However many conceptual questions surround this theorem. In what
sense can the occupation numbers vary as the model KS state is by definition
an Independent Particle State (IPS) with integer occupation numbers? What is
the functional relationship between the eigenvalues and the occupation number?
How does one define a KS functional that is differentiable with respect to
occupation numbers? In order to examine the meaning and content of this
theorem we review the construction of the implicitly defined Hohenberg-Kohn
(HK) and KS energy functionals and obtain an expression of Janaks theorem in
the context of constrained optimization theory. It is shown that, in fact,
Janaks theorem holds for any total energy functional defined in terms of
orbitals and occupation numbers. In particular it applies to Antisymmetrized
Geminal Power energy functionals, which also solely depend on orbitals and
occupation numbers and whose functional form is explicitly known. The
properties and structure of these orbital functionals are contrasted and
compared and subtle differences are discussed..\qquad

\end{abstract}
\date{05/20/2008}
\startpage{1}
\maketitle
\tableofcontents


\section{Introduction\label{Intro}}

Janaks theorem \cite{Janak1978} is a useful auxiliary tool in the Kuhn-Sham
(KS) \cite{KS1965} approach to Density Functional Theory (DFT), that is
helpful in analyzing the structure and meaning of the eigenvalues of the KS
Fock operator. It states that these eigenvalues are equal to the derivatives
of the exact total energy with respect to the occupation numbers of the First
Order Reduced Density Operator (FORDO) of the model KS state. However many
conceptual questions surround this theorem.

\begin{itemize}
\item In what sense can the occupation numbers vary as the model KS state is
by definition an Independent Particle State (IPS)\ with integer occupation numbers?

\item What is the functional relationship between the eigenvalues and the
occupation number?

\item How does one define a KS\ functional that is differentiable with respect
to occupation numbers?
\end{itemize}

In both KS-DFT and Hartree-Fock (HF) theory these eigenvalues correspond to
Lagrange parameters associated with the normalization constraint on the
orbitals, that constitute the Independent Particle State (IPS), and are
related to the ionization energies of electrons described by\textit{ }these
orbitals in the HF case and to the highest occupied one in the KS case. This
correspondence can be generalized to Multi Configurational Self Consistent
States (MCSCF) as described, for instance, in \cite{OrtizWeiner2006}.

Constrained optimization theory \cite{Hestenes75} shows, if the Hessian of the
Lagrangian is nonsingular then the values of Lagrange parameters correspond to
the sensitivity of the object function with respect to the control parameters,
i.e. to the values of the constraints. In this case the object function is the
$N$-electron energy and the control parameters are the norms of the orbitals.
Hence Janaks theorem implies that\emph{ the orbitals must be normalized to
their occupation numbers }and the energy is dependent on the occupation
numbers through its dependence on these orbitals. It is of course possible to
express the energy functional in terms of orthonormal orbitals and freely
varying occupation numbers, that comply to the basic constraints of being
positive with a maximum value of one, but this formulation does not directly
introduce the relationship between the variation of the energy functional with
respect to the occupation numbers and the eigenvalues of the Fock operator.

In general orbital functionals can be classified according to whether they are

\begin{itemize}
\item invariant to local phase changes of orbitals, i.e. are gauge invariant,
thus preserve the electron density

\item invariant to global phase changes of the orbitals i.e. to the group
$U\left(  1\right)  $ $\cite{Gilmore_{b}ook}$, e.g. Density Matrix and KS\ Functionals

\item not invariant to $U\left(  1\right)  $ e.g. correlated AGP\ functionals
\cite{OrtizWeiner2006}.
\end{itemize}

In order to answer the preceding questions and analyze the form of the energy
dependence on occupation numbers all three types of orbital functionals will
be discussed.

In this paper we will review some of the basic principles and underlying
mathematical structure of Hohenburg-Kohn (HK) \cite{HK1964} and KS DFT using
the formalism of super operators \cite{Lowdin1982} (section \ref{HKDFT}). This
will allow us to express the Levy-Lieb construction of energy functionals as
an eigenvalue problem. In section \ref{GSO} we discuss orbital based
functionals and in particular we consider KSDFT and derive Janaks Theorem in
the context of a constrained variational approach and sensitivity theory again
using the super operator formalism. The KS\ density functional, density matrix
functional and the AGP functional determine the total energy in terms of a set
of occupations numbers and one-electron orbitals. The first two have only been
defined implicitly while the AGP functional has been explicitly constructed
\cite{OrtizWeiner2006} in terms of General Spin Orbitals (GSOs), which we
review in section \ref{AGP}. In this section we compare and contrast the
properties of the derivatives of this functional with respect to occupation
numbers to the derivatives obtained indirectly through sensitivity theory for
the orbital based functional introduced in section \ref{GSO}. In section
\ref{Discussion} we conclude with a discussion and summary.

\section{Hohenberg-Kohn DFT\label{HKDFT}}

In this section we review some fundamental background needed to define the HK
density functional. We will mention important technical mathematical aspects
that need to be examined, but defer a more detailed analysis and will focus
for simplicity on the mathematically most well behaved cases.

The HK density functional is defined in terms of the kinetic energy operator,
$K,$ the coulomb operator, $V_{2},$ and the charge density $\rho$. In this
subsection we describe a construction of this functional utilizing the
formalism of super operators \cite{Lowdin1982}, that allows us to formulate
this construction as an eigenvalue problem.

\subsection{Densities and States\label{DenStat}}

Density Functional Theory is based on the existence of a universal functional,
$\mathcal{F}\left(  \rho\right)  ,$ of the charge density $\rho\left(
\boldsymbol{r}\right)  ,$ where
\begin{equation}
\rho\left(  \boldsymbol{r}\right)  =\int D^{1}\left(  \boldsymbol{r,\xi;r,\xi
}\right)  d\boldsymbol{\xi}%
\end{equation}
and $D^{1}\left(  \boldsymbol{r,\xi;r}^{\prime}\boldsymbol{,\xi}^{\prime
}\right)  $ is the integral kernel of the First Order Reduced Operator (FORDO)
$D^{1}$ (see Appendix \ref{Denmat} for details of density and reduced density
operators and matrices), which can be expanded as
\begin{equation}
D^{1}=\int D^{1}\left(  \boldsymbol{r,\xi;r}^{\prime}\boldsymbol{,\xi}%
^{\prime}\right)  \left\vert \boldsymbol{r,\xi}\right\rangle \left\langle
\boldsymbol{r}^{\prime}\boldsymbol{,\xi}^{\prime}\right\vert d\boldsymbol{r}%
d\boldsymbol{r}^{\prime}d\boldsymbol{\xi}d\boldsymbol{\xi}^{\prime}%
\end{equation}
in the basis $\left\{  \left\vert \boldsymbol{r,\xi}\right\rangle \right\}  $
formed by the tensor product of the generalized eigenvectors, $\left\{
\left\vert \boldsymbol{r}\right\rangle \right\}  ,$ of the position operator
and spin states $\left\{  \left\vert \boldsymbol{\xi}\right\rangle \right\}  $.

For a system that contains on average $N$-electrons the position density must
satisfy the conditions
\begin{subequations}
\label{DensityConst}%
\begin{align}
\int\rho\left(  \boldsymbol{r}\right)  d\boldsymbol{r}  &  =N\Rightarrow
\rho\in L_{1}\left(  \mathbb{R}^{3}\right) \label{number}\\
\rho\left(  \boldsymbol{r}\right)   &  \geq0\label{positivity}\\
\rho &  \in L_{3}\left(  \mathbb{R}^{3}\right)  . \label{KE}%
\end{align}
The constraint eq. $\left(  \ref{number}\right)  $ does not restrict the
density to be the contraction of a $N$-electron state and one could consider
incoherent and coherent mixtures of sates describing different number of
electrons. However in this article we primarily focus on $N$-electron states
and mainly consider convex sets, $\left\{  \left[  \rho\right]  _{N}\right\}
,$ of $N$-electron state density operators defined by
\end{subequations}
\begin{equation}
\left[  \rho\right]  _{N}=\left\{  \left.  D\right\vert \Xi_{1}\left(
D\right)  =\rho;D\in S_{N}\subset\mathcal{B}_{1}\left(  \mathcal{H}%
^{N}\right)  \right\}
\end{equation}
(see Appendix \ref{Denmat} for the definition of the convex set $S_{N}$ and of
$\mathcal{B}_{1}\left(  \mathcal{H}^{N}\right)  $). The constraint eq.
$\left(  \ref{KE}\right)  $ ensures that the expectation value of the kinetic
energy operator is finite \cite{Lieb1983}. The constraint eq. $\left(
\ref{positivity}\right)  $ is a necessary and sufficient condition for
representability i.e. that there exists an electronic state that contracts to
this density. The constraints eq. $\left(  \ref{number}\right)  $ and eq.
$\left(  \ref{KE}\right)  $ give that
\begin{equation}
\rho\in\mathcal{D}\subset L_{3}\left(  \mathbb{R}^{3}\right)  \cap
L_{1}\left(  \mathbb{R}^{3}\right)  ,
\end{equation}
where $\mathcal{D}$ is the convex set defined by eq. $\left(
\ref{DensityConst}\right)  $. The constraining functional $\Xi_{1}$ can be
written in terms of the set of density operators $\left\{  \left.  \Phi\left(
\boldsymbol{r}\right)  ^{\dagger}\Phi\left(  \boldsymbol{r}\right)
\right\vert \boldsymbol{r\in}\mathbb{R}^{3}\right\}  $, (see Appendix
\ref{SecQuant}) as
\begin{equation}
\Xi_{1}\left(  D\right)  \left(  \boldsymbol{r}\right)  =\operatorname{Tr}%
\left\{  \Phi\left(  \boldsymbol{r}\right)  ^{\dagger}\Phi\left(
\boldsymbol{r}\right)  D\right\}  =\rho\left(  \boldsymbol{r}\right)  .
\end{equation}
As $D$ is a positive operator it can always be expressed as a square
\begin{equation}
D=A^{2},
\end{equation}
where
\begin{equation}
A\in\mathcal{B}_{2}\left(  \mathcal{H}^{N}\right)  \text{ and }A=A^{\dagger}%
\end{equation}
and $\mathcal{B}_{2}\left(  \mathcal{H}^{N}\right)  $ is the Hilbert space of
Hilbert-Schmidt operators acting on $N$-electron space $\mathcal{H}^{N}$ (see
appendix \ref{Denmat} for the definition of $\mathcal{B}_{2}\left(
\mathcal{H}^{N}\right)  $). The operator $A$ is the square root of $D$, which
is not unique, as can be observed from the spectral expansion
\begin{equation}
A=D^{\frac{1}{2}}=\sum_{1\leq j\leq m}\pm\sqrt{\alpha_{j}}\mathcal{P}_{j},
\label{specA}%
\end{equation}
where
\begin{equation}
D=\sum_{1\leq j\leq m}\alpha_{j}\mathcal{P}_{j};\,\alpha_{j}\geq0,1\leq j\leq
m,
\end{equation}
$r_{N}=\sum\limits_{1\leq j\leq m}\nu_{j}=$ Dimension of $\mathcal{H}^{N}$,
$\left\{  \left.  \nu_{j}\right\vert 1\leq j\leq m\right\}  $ the degeneracies
of the eigenvalues $\left\{  \left.  \alpha_{j}\right\vert 1\leq j\leq
m\right\}  $ and $\left\{  \left.  \mathcal{P}_{j}\right\vert 1\leq j\leq
m\right\}  $ are eigenprojectors onto the eigenspaces of $D$. If one restricts
$A$ to be positive, i.e. uses the positive square roots in eq. $\left(
\ref{specA}\right)  $, then it belongs to the cone \cite{Encyclo} of positive
Hilbert Schmidt operators and uniquely determines $D$. The domain of variation
of the $N$-electron energy functional can thus be taken to be contained in the
\emph{real} Hilbert space $\mathcal{B}_{2sa}\left(  \mathcal{H}^{N}\right)  ,$
of dimension $r_{N}^{2},$ which has an orthonormal bases
\begin{align}
\Lambda_{jk}  &  =\frac{1}{\sqrt{2}}\left\{  \left\vert \Psi_{j}\right\rangle
\left\langle \Psi_{k}\right\vert +\left\vert \Psi_{k}\right\rangle
\left\langle \Psi_{j}\right\vert \right\}  ;\quad1\leq j<k\leq r_{N}%
\nonumber\\
\Lambda_{jj}  &  =\left\vert \Psi_{j}\right\rangle \left\langle \Psi
_{j}\right\vert ;\quad1\leq j\leq r_{N}\nonumber\\
\Lambda_{jk}  &  =\frac{i}{\sqrt{2}}\left\{  \left\vert \Psi_{j}\right\rangle
\left\langle \Psi_{k}\right\vert -\left\vert \Psi_{k}\right\rangle
\left\langle \Psi_{j}\right\vert \right\}  ;\quad1\leq j>k\leq r_{N},
\end{align}
where $\left\{  \left.  \left\vert \Psi_{j}\right\rangle \right\vert 1\leq
j\leq r_{N}\right\}  $ is any complete orthonormal basis for $\mathcal{H}%
^{N},$ of course all these spaces are infinite dimensional if $r_{N}=\infty.$

\subsection{The Density Constraint\label{DenCon}}

The Levy-Leib construction \cite{Levy1979,Lieb1983} is based on defining the
value of the density functional as the minimum of an energy variation over
states constrained to have a fixed density. The density constraint can be
formulated in terms of a quadratic functional $\Xi_{2},$ acting on the space
of self adjoint Hilbert-Schmidt operators $\mathcal{B}_{2sa}\left(
\mathcal{H}^{N}\right)  $, based on the linear functional $\Xi_{1},$ where
\begin{align}
\Xi_{2}  &  :\mathcal{B}_{2sa}\left(  \mathcal{H}^{N}\right)  \rightarrow
\mathcal{D}\nonumber\\
\Xi_{2}\left(  A\right)   &  =\operatorname{Tr}\left\{  A\Phi^{\dagger}\Phi
A\right\}  =\left(  A\left\vert \Phi^{\dagger}\Phi A\right.  \right)  =\rho,
\label{QuadCon}%
\end{align}
where the real symmetric inner product $\left(  \cdot\left\vert \cdot\right.
\right)  $ in $\mathcal{B}_{2sa}\left(  \mathcal{H}^{N}\right)  $ is defined
as
\begin{equation}
\left(  X\left\vert Y\right.  \right)  =\operatorname{Tr}\left\{  XY\right\}
. \label{Opinner}%
\end{equation}
The quadratic constraint expressed in eq. $\left(  \ref{QuadCon}\right)  $ has
the form of an expectation value of a super operator
\begin{equation}
\widehat{T}:\mathcal{B}_{2sa}\left(  \mathcal{H}^{N}\right)  \rightarrow
\mathcal{B}_{2sa}\left(  \mathcal{H}^{N}\right)  ,
\end{equation}
which can be defined in terms of a quadratic form based on the inner product
eq. $\left(  \ref{Opinner}\right)  $ (Appendix \ref{super}). The action of
such super operators, $\widehat{T},$ can be extended by
\begin{equation}
\widehat{T}\left\vert Y\right)  =\left\vert \left[  TY\right]  _{+}\right)
,\quad\left\vert Y\right)  \in\mathcal{B}_{2}\left(  \mathcal{H}^{N}\right)
\end{equation}
to the larger Hilbert-Schmidt space $\mathcal{B}_{2}\left(  \mathcal{H}%
^{N}\right)  ,$ which can be expressed as the direct sum
\begin{equation}
\mathcal{B}_{2}\left(  \mathcal{H}^{N}\right)  =\mathcal{B}_{2sa}\left(
\mathcal{H}^{N}\right)  \oplus\mathcal{B}_{2asa}\left(  \mathcal{H}%
^{N}\right)
\end{equation}
of the orthogonal subspaces formed by the self adjoint, $\mathcal{B}%
_{2sa}\left(  \mathcal{H}^{N}\right)  $, and anti self adjoint, $\mathcal{B}%
_{2asa}\left(  \mathcal{H}^{N}\right)  ,$ operators, where
\begin{equation}
T:\operatorname{Dom}\left\{  T\right\}  \subseteq\mathcal{H}^{N}%
\rightarrow\mathcal{H}^{N}%
\end{equation}
and $\left\vert \left[  TY\right]  _{+}\right)  $ is the orthogonal projection
of $TY$ onto $\mathcal{B}_{2sa}\left(  \mathcal{H}^{N}\right)  $ i.e. the self
adjoint part of $TY,$ given by
\[
\left[  TY\right]  _{+}=\frac{1}{2}\left\{  TY+YT\right\}  =\frac{1}{2}\left[
T,Y\right]  _{+}.
\]
This extension is sufficiently general for our purposes as we only consider
expressions of the form $\operatorname{Tr}\left\{  XTY\right\}  $ where
$X\in\mathcal{B}_{2sa}\left(  \mathcal{H}^{N}\right)  $ and thus is orthogonal
to the subspace of anti self adjoint operators giving
\begin{equation}
\operatorname{Tr}\left\{  XTY\right\}  =\operatorname{Tr}\left\{  X\left[
TY\right]  _{+}\right\}  .
\end{equation}
The constraint eq. $\left(  \ref{QuadCon}\right)  $ can be formulated in terms
of the \emph{generalized} \emph{vector} super operator $\widehat{\Phi
^{\dagger}\Phi}$ as%
\begin{equation}
\Xi_{2}\left(  A\right)  =\left(  A\left\vert \widehat{\Phi^{\dagger}\Phi
}A\right.  \right)  =\rho,
\end{equation}
which has super operator components, $\widehat{\Phi\left(  \boldsymbol{r}%
\right)  ^{\dagger}\Phi\left(  \boldsymbol{r}\right)  }$.

One can form the differential, $d\Xi_{2}\left(  A\right)  ,$ \cite{Encyclo}
\begin{equation}
d\Xi_{2}\left(  A\right)  :\mathcal{B}_{2sa}\left(  \mathcal{H}^{N}\right)
\rightarrow L_{1}\left(  \mathbb{R}^{3}\right)  \cap L_{3}\left(
\mathbb{R}^{3}\right)
\end{equation}
at any point $A\in\mathcal{B}_{2sa}\left(  \mathcal{H}^{N}\right)  ,$ which
has components
\begin{equation}
\left[  d\Xi_{2}\left(  A\right)  \right]  \left(  \boldsymbol{r}\right)
=\left(  \Phi\left(  \boldsymbol{r}\right)  ^{\dagger}\Phi\left(
\boldsymbol{r}\right)  A\right\vert .
\end{equation}
The set of operators $\left\{  \left.  \left\vert \Phi\left(  \boldsymbol{r}%
\right)  ^{\dagger}\Phi\left(  \boldsymbol{r}\right)  A\right)  \right\vert
\boldsymbol{r\in}\mathbb{R}^{3}\right\}  $ forms a generalized basis $\left(
\text{Appendix \ref{Bases}}\right)  $ for the normal spaces of the surfaces
$\Xi_{2}\left(  A\right)  =\rho$ for all points $A$ $\left(  \text{Appendix
\ref{Regular}}\right)  $ which shows that $d\Xi_{2}\left(  A\right)  $ is
surjective \cite{Encyclo} and thus the constraint $\Xi_{2}\left(  A\right)
=\rho$ is regular \cite{Hestenes75} everywhere.

\subsection{Hohenberg-Kohn Density Functional Construction}

Even though the explicit form of the Hohenberg-Kohn functional $\mathcal{F}%
_{HK}\left(  \rho\right)  $ is not known one can define it implicitly
\cite{Levy1979,Lieb1983} by
\begin{equation}
\mathcal{F}_{HK}\left(  \rho\right)  =\operatorname{Tr}\left\{  \left(
K+V_{2}\right)  A_{\ast}\left(  \rho\right)  ^{2}\right\}  =E_{U}\left(
A_{\ast}\left(  \rho\right)  \right)
\end{equation}
in terms of the minima, $E_{U}\left(  A_{\ast}\left(  \rho\right)  \right)  ,$
of the universal energy functional $E_{U}\left(  A\right)  ,$ that are
solutions of the constrained quadratic optimization problem
\begin{subequations}
\label{HKQ}%
\begin{align}
\min_{A\in\mathcal{B}_{2sa}\left(  \mathcal{H}^{N}\right)  }\operatorname{Tr}%
\left\{  A\overset{}{\left(  K+V_{2}\right)  A}\right\}   &  =\left(  A_{\ast
}\left(  \rho\right)  \left\vert \left(  \widehat{K}+\widehat{V_{2}}\right)
A_{\ast}\left(  \rho\right)  \right.  \right)  =E_{U}\left(  A_{\ast}\left(
\rho\right)  \right) \label{HKQa}\\
\text{subject to }\Xi_{2}\left(  A\right)   &  =\left(  A\left\vert
\widehat{\Phi^{\dagger}\Phi}A\right.  \right)  =\rho, \label{HKQb}%
\end{align}
when the control densities $\rho\in\mathcal{D}.$ This choice of the density
$\rho$ coupled with insisting that $A\in\mathcal{B}_{2sa}\left(
\mathcal{H}^{N}\right)  $ ensures that $A$ is normalized (Appendix
\ref{Normalization}$),$ i.e.
\end{subequations}
\begin{equation}
\left(  A\left\vert A\right.  \right)  =1.
\end{equation}
As this is a convex problem a local minimum is a global minimum, which,
however, may be degenerate. The Lagrangian associated with this problem eq.
$\left(  \ref{HKQ}\right)  $ is
\begin{equation}
\mathcal{L}\left(  A,\lambda,\rho\right)  =E_{U}\left(  A\right)
-\left\langle \left.  \lambda\right\vert \left\{  \Xi_{2}\left(  A\right)
-\rho\right\}  \right\rangle , \label{Lar}%
\end{equation}
where the universal energy is
\begin{equation}
E_{U}\left(  A\right)  =\operatorname{Tr}\left\{  \left(  K+V_{2}\right)
A^{2}\right\}  =\left(  A\left\vert \left(  \widehat{K}+\widehat{V_{2}%
}\right)  A\right.  \right)  .
\end{equation}
The Lagrange multiplier term can be expanded as
\begin{equation}
\left\langle \left.  \lambda\right\vert \left\{  \Xi_{2}\left(  A\right)
-\rho\right\}  \right\rangle =\int\lambda\left(  \boldsymbol{r}\right)
\left\{  \Xi_{2}\left(  A\right)  \left(  \boldsymbol{r}\right)  -\rho\left(
\boldsymbol{r}\right)  \right\}  d\boldsymbol{r},\boldsymbol{\,}\lambda\in
L_{\infty}\left(  \mathbb{R}^{3}\right)  \oplus L_{\frac{3}{2}}\left(
\mathbb{R}^{3}\right)  ,\rho\in\mathcal{D},
\end{equation}
where $\left\langle \cdot|\cdot\right\rangle $ is the scalar product between
the dual spaces $\left\langle L_{\infty}\left(  \mathbb{R}^{3}\right)  \oplus
L_{\frac{3}{2}}\left(  \mathbb{R}^{3}\right)  ,L_{1}\left(  \mathbb{R}%
^{3}\right)  \cap L_{3}\left(  \mathbb{R}^{3}\right)  \right\rangle $ i.e.
\begin{equation}
\left\langle \cdot|\cdot\right\rangle :\left(  L_{\infty}\left(
\mathbb{R}^{3}\right)  \oplus L_{\frac{3}{2}}\left(  \mathbb{R}^{3}\right)
\right)  \times\left(  L_{1}\left(  \mathbb{R}^{3}\right)  \cap L_{3}\left(
\mathbb{R}^{3}\right)  \right)  \rightarrow\mathbb{R}.
\end{equation}
This constraint term in the Lagrangian eq. $\left(  \ref{Lar}\right)  $ can
also be expressed in the form
\begin{equation}
\left\langle \left.  \lambda\right\vert \left\{  \Xi_{2}\left(  A\right)
-\rho\right\}  \right\rangle =\left\langle \lambda\left\vert \rho\right.
\right\rangle -\left(  A\left\vert \widehat{\Omega\left(  \lambda\right)
}\right.  A\right)  =\left(  A\left\vert \left\{  \left\langle \lambda
\left\vert \rho\right.  \right\rangle \hat{I}-\widehat{\Omega\left(
\lambda\right)  }\right\}  \right.  A\right)  ,
\end{equation}
where $\Omega\left(  \lambda\right)  $ is a local one body potential defined
as
\begin{equation}
\Omega\left(  \lambda\right)  =\int\lambda\left(  \boldsymbol{r}\right)
\boldsymbol{\Phi}\left(  \boldsymbol{r}\right)  ^{\dagger}\boldsymbol{\Phi
}\left(  \boldsymbol{r}\right)  d\boldsymbol{r}%
\end{equation}
and the associated super operator expectation values are given by
\begin{equation}
\left(  A\left\vert \widehat{\Omega\left(  \lambda\right)  }\right.  A\right)
=\int\lambda\left(  \boldsymbol{r}\right)  \Xi_{2}\left(  A\right)  \left(
\boldsymbol{r}\right)  d\boldsymbol{r.}%
\end{equation}
The Lagrangian for a fixed $N$-particle density can then be recast as
\begin{equation}
\mathcal{L}\left(  A,\lambda,\rho\right)  =\left(  A\left\vert \left(
\widehat{K}+\widehat{V_{2}}-\widehat{\Omega\left(  \lambda\right)  }\right)
A\right.  \right)  +\left\langle \lambda\left\vert \rho\right.  \right\rangle
.
\end{equation}


\subsection{Optimization Conditions\label{OpCon}}

Solutions , $A_{\ast},$ of the constrained quadratic optimization eq. $\left(
\ref{HKQ}\right)  $ correspond to unconstrained local saddle points $\left(
A_{\ast},\lambda_{\ast}\right)  $ \cite{Hestenes75} of the Lagrangian
$\mathcal{L}\left(  A,\lambda,\rho\right)  $
\begin{equation}
\mathcal{L}\left(  A_{\ast},\lambda_{\ast},\rho\right)  =\max_{\lambda}%
\min_{A}\mathcal{L}\left(  A,\lambda,\rho\right)
\end{equation}
that satisfy the necessary condition
\begin{equation}
d\mathcal{L}\left(  A_{\ast},\lambda_{\ast},\rho\right)  =0,
\end{equation}
where the linear map
\begin{equation}
d\mathcal{L}\left(  A,\lambda,\rho\right)  :\mathcal{B}_{2sa}\left(
\mathcal{H}^{N}\right)  \times\left(  L_{\infty}\left(  \mathbb{R}^{3}\right)
\oplus L_{\frac{3}{2}}\left(  \mathbb{R}^{3}\right)  \right)  \rightarrow
\mathbb{R}%
\end{equation}
is given by%
\begin{equation}
d\mathcal{L}\left(  A,\lambda,\rho\right)  \left(  h,\varphi\right)
=\left\langle \varphi\left\vert \rho\right.  \right\rangle -\left(
A\left\vert \widehat{\Omega\left(  \varphi\right)  }\right.  A\right)
+2\left(  h\left\vert \left\{  \widehat{K}+\widehat{V_{2}}-\widehat
{\Omega\left(  \lambda\right)  }\right\}  \right.  A\right)  . \label{LagDiff}%
\end{equation}
As the constraint is regular the Lagrange functional $\lambda_{\ast}$ is
unique \cite{Hestenes75}, but there may be degenerate $A_{\ast}$'s. The
differential eq. $\left(  \ref{LagDiff}\right)  $ can be partitioned into
$\left(  d_{A}\mathcal{L}\left(  A,\lambda,\rho\right)  ,d_{\lambda
}\mathcal{L}\left(  A,\lambda,\rho\right)  \right)  $, where
\begin{subequations}
\label{differential}%
\begin{align}
d_{A}\mathcal{L}\left(  A,\lambda,\rho\right)  \left(  \cdot\right)   &
=2\left(  \cdot\left\vert \left\{  \widehat{K}+\widehat{V_{2}}-\widehat
{\Omega\left(  \lambda\right)  }\right\}  \right.  A\right) \\
d_{\lambda}\mathcal{L}\left(  A,\lambda,\rho\right)  \left(  \cdot\right)   &
=\left\langle \cdot\left\vert \rho\right.  \right\rangle -\left(  A\left\vert
\Omega\left(  \cdot\right)  \right.  A\right)  ,
\end{align}
thus giving the stationary point conditions
\end{subequations}
\begin{subequations}
\begin{align}
\left(  h\left\vert \left\{  \widehat{K}+\widehat{V_{2}}-\widehat
{\Omega\left(  \lambda_{\ast}\right)  }\right\}  \right.  A_{\ast}\right)   &
=0\,\forall h\in\mathcal{B}_{2sa}\left(  \mathcal{H}^{N}\right)
\label{stata}\\
\left\langle \varphi\left\vert \rho\right.  \right\rangle -\left(  A_{\ast
}\left\vert \widehat{\Omega\left(  \varphi\right)  }\right.  A_{\ast}\right)
&  =0\,\forall\varphi\in L_{\infty}\left(  \mathbb{R}^{3}\right)  \oplus
L_{\frac{3}{2}}\left(  \mathbb{R}^{3}\right)  . \label{statb}%
\end{align}
If eq. $\left(  \ref{stata}\right)  $ and eq. $\left(  \ref{statb}\right)  $
are true for all operators $h$ belonging to $\mathcal{B}_{2sa}\left(
\mathcal{H}^{N}\right)  $ and all functions $\varphi$ belonging to $L_{\infty
}\left(  \mathbb{R}^{3}\right)  \oplus L_{\frac{3}{2}}\left(  \mathbb{R}%
^{3}\right)  $ these conditions are equivalent to
\end{subequations}
\begin{equation}
\left\{  \widehat{K}+\widehat{V_{2}}-\widehat{\Omega\left(  \lambda_{\ast
}\right)  }\right\}  \left\vert A_{\ast}\right)  =0 \label{Stateig}%
\end{equation}
and the constraining condition
\begin{equation}
\left(  A_{\ast}\left\vert \widehat{\Phi^{\dagger}\Phi}\right.  A_{\ast
}\right)  =\rho.
\end{equation}
(For unbounded operators $K+V_{2}-\Omega\left(  \lambda_{\ast}\right)  $ one
must be more precise, but here we will only consider the case when an
eigenvector $\left\vert A_{\ast}\right)  $ exists and thus eq. $\left(
\ref{stata}\right)  $ is true for all $h\in\mathcal{B}_{2sa}\left(
\mathcal{H}^{N}\right)  $ and $\rho$ is V-representable, more general cases
will be discussed elsewhere.

$\lceil$%
\[
\left\{  \widehat{K}+\widehat{V_{2}}-\widehat{\Omega\left(  \lambda_{\ast
}\right)  }\right\}  \left\vert A_{\ast}\right)  =0\Rightarrow
\operatorname{Tr}\left\{  h\left\{  K+V_{2}-\Omega\left(  \lambda_{\ast
}\right)  \right\}  A_{\ast}\right\}  =0\forall h
\]
$\rfloor$

The optimal value of the Lagrangian satisfies the following identities:
\begin{align}
\mathcal{L}\left(  A_{\ast},\lambda_{\ast},\rho\right)   &  =\left(  A_{\ast
}\left(  \rho\right)  \left\vert \left\{  \left\{  \widehat{K}+\widehat{V_{2}%
}-\widehat{\Omega\left(  \lambda_{\ast}\right)  }\right\}  +\left\langle
\left.  \lambda_{\ast}\left(  \rho\right)  \right\vert \rho\right\rangle
\hat{I}\right\}  \right.  A_{\ast}\left(  \rho\right)  \right)  =\left\langle
\left.  \lambda_{\ast}\left(  \rho\right)  \right\vert \rho\right\rangle
\nonumber\\
&  =\left(  A_{\ast}\left(  \rho\right)  \left\vert \left\{  \widehat
{K}+\widehat{V_{2}}\right\}  \right.  A_{\ast}\left(  \rho\right)  \right)
=E_{U}\left(  A_{\ast}\left(  \rho\right)  \right)  .
\end{align}
This shows that the optimal values of the universal energy functional, the
value of the Lagrangian and the value of the HK energy functional are all
equal at the density $\rho$ i.e.
\begin{equation}
\mathcal{F}_{HK}\left(  \rho\right)  =E_{U}\left(  A_{\ast}\left(
\rho\right)  \right)  =\left\langle \left.  \lambda_{\ast}\left(  \rho\right)
\right\vert \rho\right\rangle .
\end{equation}
Moreover $\rho$ is the density of the ground state $D_{\ast}\left(
\rho\right)  =A_{\ast}\left(  \rho\right)  ^{2}$ of the Hamiltonian
\begin{equation}
K+V_{2}-\Omega\left(  \lambda_{\ast}\left(  \rho\right)  \right)
\end{equation}
i.e. of an electronic system subject to an external potential $\lambda_{\ast
}\left(  \rho\right)  $. Note that the ground state might be a ensemble state.

\subsection{Sensitivity Analysis\label{SenAnal}}

Sensitivity analysis \cite{Hestenes75} examines how the optimal solution
changes with variation of the value of the constraining functions, which can
be thought of as control parameters. In this case the $N$-electron state
energy with variations of $\rho.$ If the optimal points $\left\{  A_{\ast
},\lambda_{\ast},\rho\right\}  $ are regular (Appendix \ref{Regular}) for a
range of $\rho^{\prime}s$ then $\left\{  A_{\ast},\lambda_{\ast}\right\}  $
are in fact functions of $\rho,$ i.e.
\begin{align}
A_{\ast}  &  =A_{\ast}\left(  \rho\right) \nonumber\\
\lambda_{\ast}  &  =\lambda_{\ast}\left(  \rho\right)  ,
\end{align}
and thus single valued i.e. the ground state is non degenerate. The
sensitivity of the optimal $N$-electron energy to changes in the density
$\rho$ is then given by
\begin{equation}
\frac{\delta E_{U}\left(  A_{\ast}\left(  \rho\right)  \right)  }{\delta\rho
}=\lambda_{\ast}\left(  \rho\right)  .
\end{equation}
By construction
\begin{equation}
\mathcal{F}_{HK}\left(  \rho\right)  =E_{U}\left(  A_{\ast}\left(
\rho\right)  \right)  ,
\end{equation}
which shows
\begin{equation}
\frac{\delta\mathcal{F}_{HK}\left(  \rho\right)  }{\delta\rho}=\lambda_{\ast
}\left(  \rho\right)  .
\end{equation}


\subsection{Decomposition of the Hohenburg-Kohn Functional\label{DecompHK}}

If a minimizer $D_{\ast}\left(  \rho\right)  $ exists and is unique and then
one can define kinetic energy, coulomb and exchange-correlation functionals
as
\begin{align}
\mathcal{F}_{K}\left(  \rho\right)   &  =\operatorname{Tr}\left\{  KD_{\ast
}\left(  \rho\right)  \right\}  ,\nonumber\\
\mathcal{F}_{C}\left(  \rho\right)   &  =\int\frac{\rho\left(
\boldsymbol{r,\xi}\right)  \rho\left(  \boldsymbol{r}^{\prime}\boldsymbol{,\xi
}\right)  }{\left\Vert \boldsymbol{r-r}^{\prime}\right\Vert }d\boldsymbol{r}%
d\boldsymbol{r}^{\prime}d\boldsymbol{\xi},\nonumber\\
\mathcal{F}_{XC}\left(  \rho\right)   &  =\operatorname{Tr}\left\{
V_{2}D_{\ast}\left(  \rho\right)  \right\}  -\mathcal{F}_{C}\left(
\rho\right)  ,
\end{align}
producing the decomposition
\begin{equation}
\mathcal{F}_{HK}\left(  \rho\right)  =\mathcal{F}_{K}\left(  \rho\right)
+\mathcal{F}_{C}\left(  \rho\right)  +\mathcal{F}_{XC}\left(  \rho\right)  .
\end{equation}
The concept of a unique exchange-correlation functional is central to the
theory, but the existence of non isolated degenerate $D_{\ast}\left(
\rho\right)  ^{\prime}$s could give different $\mathcal{F}_{XC}\left(
\rho\right)  ^{\prime}$s and $\mathcal{F}_{K}\left(  \rho\right)  ^{\prime}$s.
However if degeneracies do occur one can still define an exchange-correlation
functional through the imposition of further constrained minimizations.

\subsection{Molecular System Energy\label{MolSysEn}}

The total energy functional of a system that contains $N$ electrons and
experiences an environment described by a \emph{local} potential $\mathrm{v}$
is
\begin{equation}
E\left(  \rho\boldsymbol{,}\mathrm{v}\right)  =\mathcal{F}_{HK}\left(
\rho\right)  +f_{\rho}\left(  \mathrm{v}\right)  ,
\end{equation}
where the explicit effect of the environment is described by the functional
\begin{equation}
f_{\rho}\left(  \mathrm{v}\right)  =\int\mathrm{v}\left(  \boldsymbol{r}%
\right)  \rho\left(  \boldsymbol{r}\right)  d\boldsymbol{r.}%
\end{equation}
The ground state energy $E_{GS}\left(  \mathrm{v}\right)  $ and the
corresponding density $\rho_{GS}$ $\left(  \mathrm{v}\right)  $ are then
determined by the solution to the constrained optimization problem
\begin{equation}
E_{GS}\left(  \mathrm{v}\right)  =E\left(  \rho_{GS}\boldsymbol{,}%
\mathrm{v}\right)  =\min_{\rho\in\mathcal{D}}E\left(  \rho\boldsymbol{,}%
\mathrm{v}\right)  , \label{GS}%
\end{equation}
which one can make unconstrained in terms of real functions $f\in H_{1}\left(
\mathbb{R}^{3}\right)  $ \cite{Lieb1983} by setting
\begin{equation}
\rho\left(  \boldsymbol{r}\right)  =\frac{Nf\left(  \boldsymbol{r}\right)
^{2}}{\left\langle f|f\right\rangle }=\frac{N\left\langle f\left\vert
\Phi\left(  \boldsymbol{r}\right)  ^{\dagger}\Phi\left(  \boldsymbol{r}%
\right)  f\right.  \right\rangle }{\left\langle f|f\right\rangle },\,f\in
H_{1}\left(  \mathbb{R}^{3}\right)  . \label{rho_f}%
\end{equation}
The minimizer $f_{GS}$ of $E\left(  f\boldsymbol{,}\mathrm{v}\right)  $ then
satisfies
\begin{equation}
\frac{\delta E\left(  f_{GS}\boldsymbol{,}\mathrm{v}\right)  }{\delta f\left(
\boldsymbol{r}\right)  }=2f_{GS}\left(  \boldsymbol{r}\right)  \left\{
\frac{N}{\left\langle f_{GS}|f_{GS}\right\rangle }\frac{\delta E\left(
f_{GS}\boldsymbol{,}\mathrm{v}\right)  }{\delta\rho\left(  \boldsymbol{r}%
\right)  }-\int\frac{\delta E\left(  f_{GS}\boldsymbol{,}\mathrm{v}\right)
}{\delta\rho\left(  \boldsymbol{r}^{\prime}\right)  }\rho_{GS}\left(
\boldsymbol{r}^{\prime}\right)  d\boldsymbol{r}^{\prime}\right\}
\overset{a.e.}{=}0
\end{equation}
which implies that
\begin{equation}
f_{GS}\left(  \boldsymbol{r}\right)  \overset{a.e.}{=}0\Leftrightarrow
\rho_{GS}\left(  \boldsymbol{r}\right)  \overset{a.e.}{=}0
\end{equation}
or
\begin{equation}
\frac{\delta E\left(  f_{GS}\boldsymbol{,}\mathrm{v}\right)  }{\delta
\rho\left(  \boldsymbol{r}\right)  }\overset{a.e.}{=}\alpha\left(  \rho
_{GS}\right)  , \label{fmin}%
\end{equation}
where%
\begin{equation}
\alpha\left(  \rho_{GS}\right)  =\frac{\left\langle f_{GS}|f_{GS}\right\rangle
}{N}\int\frac{\delta E\left(  f_{GS}\boldsymbol{,}\mathrm{v}\right)  }%
{\delta\rho\left(  \boldsymbol{r}^{\prime}\right)  }\rho_{GS}\left(
\boldsymbol{r}^{\prime}\right)  d\boldsymbol{r}^{\prime}. \label{alpha}%
\end{equation}
Hence
\begin{align}
\alpha\left(  \rho_{GS}\right)   &  =\frac{\left\langle f_{GS}|f_{GS}%
\right\rangle }{N}\alpha\left(  \rho_{GS}\right)  \int\rho_{GS}\left(
\boldsymbol{r}^{\prime}\right)  d\boldsymbol{r}^{\prime}=\frac{\left\langle
f_{GS}|f_{GS}\right\rangle }{N}\alpha\left(  \rho_{GS}\right) \\
&  \Rightarrow\alpha\left(  \rho_{GS}\right)  =0\text{ or }f_{GS}\left(
\boldsymbol{r}\right)  \overset{a.e.}{=}0.
\end{align}
So as $f_{GS}\left(  \boldsymbol{r}\right)  \overset{a.e.}{\neq}0$ we obtain
the well known result
\begin{equation}
\frac{\delta E\left(  f_{GS}\boldsymbol{,}\mathrm{v}\right)  }{\delta
\rho\left(  \boldsymbol{r}\right)  }\overset{a.e.}{=}0\Leftrightarrow
\frac{\delta\mathcal{F}_{HK}\left(  \rho_{GS}\right)  }{\delta\rho\left(
\boldsymbol{r}\right)  }\overset{a.e.}{=}-\mathrm{v}\left(  \boldsymbol{r}%
\right)  .
\end{equation}


\section{Orbital Variables \label{GSO}}

There are two standard approaches to formulating the total energy in terms of
orbitals and occupation numbers, one is to consider functionals of FORDO's,
which leads to DMFT, the other is to factor the dependence through the
density, which leads to KS-DFT and is based on the HK functional. However as
the explicit forms of these functionals are not known one must rely on the
implicit definitions via constrained optimizations, though in practice
approximate explicit forms are utilized. We have recently
\cite{OrtizWeiner2006} introduced a third approach which is defined on states
generated by the submanifold of $2M$-electron AGP states. The total energy
functional, which is explicitly given, \emph{is more general }than a density
matrix functional while still expressing the total energy in terms of orbitals
and occupation numbers.

The central focus of this article is on functionals of the density that
contain no explicit reference to the spin properties of the system. However
one could widen this focus by considering spin and general spin densities that
are described by vectors $\boldsymbol{\rho}\left(  \boldsymbol{r}\right)
=\left(  \rho_{\alpha}\left(  \boldsymbol{r}\right)  ,\rho_{\beta}\left(
\boldsymbol{r}\right)  \right)  $ and $\boldsymbol{\rho}\left(  \boldsymbol{r}%
\right)  =\left(  \rho_{\alpha}\left(  \boldsymbol{r}\right)  ,\rho
_{\alpha\beta}\left(  \boldsymbol{r}\right)  ,\rho_{\beta}\left(
\boldsymbol{r}\right)  \right)  $ respectively and modify the preceding
analysis in the appropriate fashion (Appendix \ref{Spin}).

\subsection{Correspondence between Orbitals and Densities\label{Orb_Den}}

Any density $\rho$ can expressed in terms of the Natural Orbitals (NO's)
$\left\{  \psi_{j}\right\}  $ and occupation numbers $\left\{  N_{j}\right\}
$ of the FORDO $D_{\ast}^{1}\left(  \rho\right)  $ (Appendix \ref{Denmat}) as%

\begin{equation}
\rho\left(  \boldsymbol{r}\right)  =\left[  D_{\ast}^{1}\left(  \rho\right)
\right]  \left(  \boldsymbol{r,r}\right)  =\sum_{1\leq j\leq r}N_{j}\left\vert
\psi_{j}\left(  \boldsymbol{r}\right)  \right\vert ^{2}, \label{NGSOexpan}%
\end{equation}
but many other sets of orbitals also produce the same density, in particular
Harimanns \cite{Harriman81} theorem ensures that there are always expansions
of the form
\begin{equation}
\rho\left(  \boldsymbol{r}\right)  =\left[  D_{\ast}^{1}\left(  \rho\right)
\right]  \left(  \boldsymbol{r,r}\right)  =\sum_{1\leq j\leq N}\chi_{j}\left(
\boldsymbol{r}\right)  ^{2} \label{Genexp}%
\end{equation}
in terms of real orthonormal square summable functions $\left\{  \left.
\chi_{j}\right\vert 1\leq j\leq N\right\}  $.

In general one can express the density in terms of orthogonal square
integrable functions as
\begin{equation}
\rho\left(  \boldsymbol{r}\right)  =\sum_{1\leq j\leq\nu}\left\vert
\varphi_{j}\left(  \boldsymbol{r}\right)  \right\vert ^{2}, \label{OrthogOrb}%
\end{equation}
where
\begin{equation}
\left\langle \left.  \varphi_{j}\right\vert \varphi_{k}\right\rangle
=N_{j}\delta_{jk};0<N_{j}\leq N,1\leq j\leq\nu;N_{j}=0,\nu+1\leq j\leq r
\end{equation}
and $\nu\equiv\nu\left(  \boldsymbol{\varphi}\right)  $ is called the rank of
the expansion. If $\nu\geq N$ then the functions $\boldsymbol{\varphi}=\left(
\varphi_{1},\cdots,\varphi_{r}\right)  $ can be interpretted as orbitals. One
can also expand the density $\rho$ in terms of orthonormal functions
$\boldsymbol{\tilde{\varphi}}$ and weights $\boldsymbol{N}=\left(
N_{1},\cdots,N_{r}\right)  $ as
\begin{equation}
\rho\left(  \boldsymbol{r}\right)  =\sum_{1\leq j\leq\nu}N_{j}\left\vert
\tilde{\varphi}_{j}\left(  \boldsymbol{r}\right)  \right\vert ^{2}.
\label{orthonormal}%
\end{equation}
Again if $\nu\geq N$ then $\boldsymbol{\tilde{\varphi}}$ and $\boldsymbol{N}$
can be interpretted as normalized orbitals and occupation numbers
respectively. One can make the relationship expressed eq. $\left(
\ref{OrthogOrb}\right)  $ and eq. $\left(  \ref{orthonormal}\right)  $
explicit through the introduction of two functionals $\Gamma_{1}$ and
$\Gamma_{2}.$ The first functional $\Gamma_{1}$ is defined by
\begin{equation}
\Gamma_{1}:\mathfrak{S}_{1}\subset\prod\limits_{1\leq j\leq r}\mathcal{H}%
^{1}\rightarrow L_{1+}\left(  \mathbb{R}^{3}\right)  \subset L_{1}\left(
\mathbb{R}^{3}\right)  ,
\end{equation}
where
\begin{subequations}
\label{GammaAll}%
\begin{align}
\rho &  =\Gamma_{1}\left(  \boldsymbol{\varphi}\right)  ,\\
\left[  \Gamma_{1}\left(  \boldsymbol{\varphi}\right)  \right]  \left(
\boldsymbol{r}\right)   &  =\sum_{1\leq j\leq r}\left\vert \varphi_{j}\left(
\boldsymbol{r}\right)  \right\vert ^{2}=\rho\left(  \boldsymbol{r}\right)
\label{Gamma}%
\end{align}
and%
\end{subequations}
\begin{equation}
\mathfrak{S}_{1}=\left\{  \boldsymbol{\varphi}\left\vert \varphi_{j}%
,\varphi_{k}\in\mathcal{H}^{1},\left\langle \left.  \varphi_{j}\right\vert
\varphi_{k}\right\rangle =N_{j}\delta_{jk},\,\boldsymbol{N}\in\mathfrak{N}%
\left(  N\right)  \right.  \right\}  , \label{set1}%
\end{equation}
where
\begin{equation}
\mathfrak{N}\left(  N\right)  =\left\{  \boldsymbol{N}\left\vert 0\leq
N_{j}\leq N,1\leq j\leq r;\sum_{1\leq j\leq r}N_{j}=N\right.  \right\}  .
\end{equation}
and as many as $r-1$ of the component functions, $\varphi_{j},$ may be
identically zero. The set $\mathfrak{S}_{1}$ includes all possible expansions
of $\rho$ even the one term case considered in eq. $\left(  \ref{rho_f}%
\right)  $. The second functional $\Gamma_{2}$ is defined by
\begin{equation}
\Gamma_{2}:\mathfrak{S}_{2}\subset\prod\limits_{1\leq j\leq r}\mathcal{H}%
^{1}\times\mathbb{R}^{r}\rightarrow L_{1+}\left(  \mathbb{R}^{3}\right)
\subset L_{1}\left(  \mathbb{R}^{3}\right)  ,
\end{equation}
where
\begin{equation}
\mathfrak{S}_{2}=\left\{  \left(  \boldsymbol{\tilde{\varphi},N}\right)
\left\vert \left\vert \tilde{\varphi}_{j}\right\rangle ,\left\vert
\tilde{\varphi}_{k}\right\rangle \in\mathcal{H}^{1},\left\langle \left.
\varphi_{j}\right\vert \varphi_{k}\right\rangle =\delta_{jk},\,1\leq j\leq
r;\quad\boldsymbol{N}\in\mathfrak{N}\left(  N\right)  \right.  \right\}  .
\label{set2}%
\end{equation}
Clearly it is always possible to find $\boldsymbol{\varphi}$ and $\left(
\boldsymbol{\tilde{\varphi},N}\right)  $ such that
\begin{equation}
\rho=\Gamma_{1}\left(  \boldsymbol{\varphi}\right)  =\Gamma_{2}\left(
\boldsymbol{\tilde{\varphi},N}\right)
\end{equation}
and if the number of non zero components in $\boldsymbol{\varphi}$ and
$\boldsymbol{N}$ are greater than or equal to $N$ then $\boldsymbol{\varphi}$
and $\left\{  \boldsymbol{\tilde{\varphi},N}\right\}  $ determine the same
FORDO $D^{1}$ i.e.
\begin{equation}
D^{1}\left(  \boldsymbol{\varphi}\right)  =D^{1}\left(  \boldsymbol{\tilde
{\varphi},N}\right)  =\sum_{1\leq j\leq r}\left\vert \varphi_{j}\right\rangle
\left\langle \varphi_{j}\right\vert =\sum_{1\leq j\leq r}\left\vert
\tilde{\varphi}_{j}\right\rangle \left\langle \tilde{\varphi}_{j}\right\vert
N_{j}.
\end{equation}
The optimization wrt $\rho$ can thus be expressed in terms of
$\boldsymbol{\varphi}$ via $\Gamma_{1}$ or $\left(  \boldsymbol{\tilde
{\varphi},N}\right)  $ via $\Gamma_{2}$, However in order to use sensitivity
theory to relate the variation of the energy with respect to occupation
numbers to the eigenvalues of the Kohn-Sham operator we focus on $\Gamma_{1}.$

The functional $\Gamma_{1}$ is not one to one and is constant on the
equivalence class $\mathfrak{S}\left(  \rho\right)  $ of functions%
\begin{equation}
\mathfrak{S}\left(  \rho\right)  =\bigcup\limits_{1\leq\nu\leq r}%
\mathfrak{S}\left(  \nu,\rho\right)  ,
\end{equation}
where the equivalence classes of constant rank, $\mathfrak{S}\left(  \nu
,\rho\right)  ,$ are defined by
\begin{equation}
\mathfrak{S}\left(  \nu,\rho\right)  =\left\{  \left.  \boldsymbol{\varphi
}\right\vert \nu=\nu\left(  \boldsymbol{\varphi}\right)  ,\Gamma_{1}\left(
\boldsymbol{\varphi}\right)  =\rho\right\}  .
\end{equation}
In the following we will limit consideration to the case where $\nu\left(
\boldsymbol{\varphi}\right)  \geq N$ i.e. to orbitals. Each preimage set
$\Gamma_{1}^{-1}\left(  \rho\right)  =\mathfrak{S}\left(  \rho\right)
\subset\mathfrak{S}_{1}$ contains, in general, both single sets of orthogonal
orbitals normalized to different $\boldsymbol{N}\in\mathfrak{N}\left(
N\right)  $ and different sets of orthogonal functions normalized to the same
$\boldsymbol{N}\in\mathfrak{N}\left(  N\right)  ,$ which corresponds to the
fact that a given density can be expanded in many ways in terms of occupation
numbers and orbitals.

\subsection{Orbital functionals\label{GSOFunc}}

One can express the total energy via a universal orbital functional,
$\mathcal{F}_{O}$, whose value is defined by the HK energy functional, as
\begin{equation}
\mathcal{F}_{O}\left(  \boldsymbol{\varphi}\right)  =\mathcal{F}_{HK}\left(
\Gamma_{1}\left(  \boldsymbol{\varphi}\right)  \right)  =\mathcal{F}%
_{HK}\left(  \rho\right)  . \label{FO}%
\end{equation}
The functional $\mathcal{F}_{O}$ can be decomposed as
\begin{equation}
\mathcal{F}_{O}\left(  \boldsymbol{\varphi}\right)  =\mathcal{F}_{T}\left(
\boldsymbol{\varphi}\right)  +\mathcal{F}_{C}\left(  \boldsymbol{\varphi
}\right)  +\mathcal{F}_{OXC}\left(  \boldsymbol{\varphi}\right)  , \label{KS1}%
\end{equation}
where the orbital kinetic energy functional $\mathcal{F}_{T}$ is given by
\begin{equation}
\mathcal{F}_{T}\left(  \boldsymbol{\varphi}\right)  =\sum_{1\leq j\leq
r}\left\langle \varphi_{j}\left\vert K\varphi_{j}\right.  \right\rangle ,
\end{equation}
the orbital exchange-corellation functional $\mathcal{F}_{OXC}$ (which depends
on $\mathcal{F}_{T}\boldsymbol{)}$, is defined as
\begin{equation}
\mathcal{F}_{OXC}\left(  \boldsymbol{\varphi}\right)  =\mathcal{F}_{XC}\left(
\Gamma_{1}\left(  \boldsymbol{\varphi}\right)  \right)  +\mathcal{F}%
_{K}\left(  \Gamma_{1}\left(  \boldsymbol{\varphi}\right)  \right)
-\mathcal{F}_{T}\left(  \boldsymbol{\varphi}\right)  , \label{KS2}%
\end{equation}
and one recalls the definitions
\begin{equation}
\mathcal{F}_{XC}\left(  \Gamma_{1}\left(  \boldsymbol{\varphi}\right)
\right)  =\operatorname{Tr}\left\{  V_{2}D_{\ast}\left(  \Gamma_{1}\left(
\boldsymbol{\varphi}\right)  \right)  \right\}  -\mathcal{F}_{C}\left(
\Gamma_{1}\left(  \boldsymbol{\varphi}\right)  \right)  \label{KS3}%
\end{equation}
and
\begin{equation}
\mathcal{F}_{K}\left(  \Gamma_{1}\left(  \boldsymbol{\varphi}\right)  \right)
=\operatorname{Tr}\left\{  KD_{\ast}\left(  \Gamma_{1}\left(
\boldsymbol{\varphi}\right)  \right)  \right\}  . \label{KS4}%
\end{equation}
One should note that

\begin{itemize}
\item for fixed values of $\rho$ the HK functionals $\mathcal{F}_{XC}\left(
\Gamma_{1}\left(  \boldsymbol{\varphi}\right)  \right)  $ and $\mathcal{F}%
_{K}\left(  \Gamma_{1}\left(  \boldsymbol{\varphi}\right)  \right)  $ in eqs.
$\left(  \ref{KS3}\right)  $ and $\left(  \ref{KS4}\right)  $ are
\emph{constant} for all $\boldsymbol{\varphi}$ that satisfy $\rho=\Gamma
_{1}\left(  \boldsymbol{\varphi}\right)  $, while the orbital functionals
$\mathcal{F}_{T}\left(  \boldsymbol{\varphi}\right)  $ and $\mathcal{F}%
_{OXC}\left(  \boldsymbol{\varphi}\right)  $ can have \emph{different} values
for values of $\boldsymbol{\varphi}$ that correspond to a fixed values of
$\rho$

\item the total energy functionals $\mathcal{F}_{O}\left(  \boldsymbol{\varphi
}\right)  ,\mathcal{\ }$and $\mathcal{F}_{HK}\left(  \rho\right)  $ have the
same value when $\rho=$ $\Gamma_{1}\left(  \boldsymbol{\varphi}\right)
\equiv\rho\left(  \boldsymbol{\varphi}\right)  $

\item the composition in eq. $\left(  \ref{FO}\right)  $ defines a class of
orbital functionals that have a particular form. One could define orbital
functionals without reference to densities (as done in section \ref{AGPSect})
and their form and properties would be quite different.
\end{itemize}

The minimization problem that determines the ground state energy and density
of a system characterized by a local potential $\mathrm{v}$ is
\begin{equation}
E_{GS}\left(  \mathrm{v}\right)  =\mathcal{E}\left(  \Gamma_{1}\left(
\boldsymbol{\varphi}_{GS}\right)  \boldsymbol{,}\mathrm{v}\right)
=\min_{\boldsymbol{\varphi\in}\mathfrak{S}_{1}}\mathcal{E}_{O}\left(
\boldsymbol{\varphi,}\mathrm{v}\right)  , \label{KSmin}%
\end{equation}
where
\begin{equation}
\mathcal{E}_{O}\left(  \boldsymbol{\varphi,}\mathrm{v}\right)  =\mathcal{F}%
_{O}\left(  \boldsymbol{\varphi}\right)  +f_{\Gamma_{1}\left(
\boldsymbol{\varphi}\right)  }\left(  \mathrm{v}\right)
\end{equation}
and the feasible set $\mathfrak{S}_{1}$ is defined in eq. $\left(
\ref{set1}\right)  $.

The orbital minimization problem eq. $\left(  \ref{KSmin}\right)  $ can be
decomposed into two sequential optimizations:

(a)\ first one varies the orbitals keeping their normalization fixed and
determines the saddle point $\mathcal{L}_{\boldsymbol{\varphi}}\left(
\boldsymbol{\varphi}\left(  \boldsymbol{N}\right)  ,\boldsymbol{\varepsilon
}\left(  \boldsymbol{N}\right)  ,\mu\left(  N\right)  ,\mathrm{v}\right)  $ of
the Lagrangian
\begin{equation}
\mathcal{L}_{\boldsymbol{\varphi}}\left(  \boldsymbol{\varphi}%
,\boldsymbol{\varepsilon,N,}\mu\boldsymbol{,}\mathrm{v}\right)  =\mathcal{E}%
_{O}\left(  \boldsymbol{\varphi,}\mathrm{v}\right)  -\sum_{1\leq j,k\leq
r}\varepsilon_{jk}\left\{  \overset{\,}{\left\langle \left.  \varphi
_{j}\right\vert \varphi_{k}\right\rangle -N_{j}\delta_{jk}}\right\}
+\mu\left\{  \sum_{\,1\leq j\leq r}\left\langle \left.  \varphi_{j}\right\vert
\varphi_{j}\right\rangle -N\right\}  \label{OrbLag}%
\end{equation}
i.e.
\begin{equation}
\mathcal{L}_{\boldsymbol{\varphi}}\left(  \boldsymbol{\varphi}\left(
\boldsymbol{N}\right)  ,\boldsymbol{\varepsilon}\left(  \boldsymbol{N}\right)
,\mu\left(  N\right)  ,\mathrm{v}\right)  =\max_{\boldsymbol{\varepsilon,}\mu
}\min_{\boldsymbol{\varphi}}\mathcal{L}_{\boldsymbol{\varphi}}\left(
\boldsymbol{\varphi},\boldsymbol{\varepsilon},\mu\boldsymbol{,N,}%
\mathrm{v}\right)  , \label{KSM1}%
\end{equation}
which leads to optimal values $\mathcal{E}_{O}\left(  \boldsymbol{\varphi
}\left(  \boldsymbol{N}\right)  \boldsymbol{,}\mathrm{v}\right)  =$
$\mathcal{L}_{\boldsymbol{\varphi}}\left(  \boldsymbol{\varphi}\left(
\boldsymbol{N}\right)  ,\boldsymbol{\varepsilon}\left(  \boldsymbol{N}\right)
,\mu\left(  N\right)  \boldsymbol{,}\mathrm{v}\right)  $ of the orbital energy
functional for fixed values of the control parameters $\boldsymbol{N}$. The
particle number constraint $\sum_{\,1\leq j\leq r}\left\langle \left.
\varphi_{j}\right\vert \varphi_{j}\right\rangle =N$ is equivalent to the
normalization constraint imposed on the density so that $\rho\in\mathcal{D}$.

The saddle points of $\mathcal{L}_{\boldsymbol{\varphi}}\left(
\boldsymbol{\varphi},\boldsymbol{\varepsilon,N,}\mu\boldsymbol{,}%
\mathrm{v}\right)  $ are characterized by the necessary conditions
\begin{subequations}
\label{KSC}%
\begin{align}
\frac{\delta\mathcal{L}_{\boldsymbol{\varphi}}\left(  \boldsymbol{\varphi
},\boldsymbol{\varepsilon,}\mu,\boldsymbol{N,}\mathrm{v}\right)  }%
{\delta\varphi_{j}^{\ast}}  &  =0;\,1\leq j\leq r,\label{KSC1}\\
\frac{\partial\mathcal{L}_{\boldsymbol{\varphi}}\left(  \boldsymbol{\varphi
},\boldsymbol{\varepsilon,}\mu,\boldsymbol{N,}\mathrm{v}\right)  }%
{\partial\varepsilon_{jk}}  &  =0;\,1\leq j,k\leq r\label{KSC2}\\
\frac{\partial\mathcal{L}_{\boldsymbol{\varphi}}\left(  \boldsymbol{\varphi
},\boldsymbol{\varepsilon,}\mu,\boldsymbol{N,}\mathrm{v}\right)  }{\partial
\mu}  &  =0 \label{KSC3}%
\end{align}
The condition eq. $\left(  \ref{KSC1}\right)  $ implies that at an optimal
point the matrix of Lagrange multipliers $\boldsymbol{\varepsilon}\left(
\boldsymbol{N}\right)  $ is hermitian i.e.
\end{subequations}
\begin{equation}
\varepsilon\left(  \boldsymbol{N}\right)  _{jk}=\varepsilon\left(
\boldsymbol{N}\right)  _{kj}^{\ast},\quad1\leq j,k\leq r
\end{equation}
and
\begin{equation}
\frac{\delta\mathcal{E}_{O}\left(  \boldsymbol{\varphi}\left(  \boldsymbol{N}%
\right)  ,\boldsymbol{\varepsilon}\left(  \boldsymbol{N}\right)
,\mathrm{v}\right)  }{\delta\varphi_{j}^{\ast}}-\sum_{1\leq k\leq r}\left(
\varepsilon_{jk}\left(  \boldsymbol{N}\right)  -\mu\left(  N\right)
\delta_{jk}\right)  \left\vert \varphi_{k}\left(  \boldsymbol{N}\right)
\right\rangle =0;\,1\leq j\leq r. \label{LagPar}%
\end{equation}
Noting that as $\rho=\Gamma_{1}\left(  \boldsymbol{\varphi}\right)  $ and
$\mathcal{E}_{O}\left(  \rho\left(  \boldsymbol{\varphi}\right)
,\mathrm{v}\right)  \equiv\mathcal{E}_{O}\left(  \boldsymbol{\varphi
},\mathrm{v}\right)  $ one can obtain
\begin{align}
\frac{\delta\mathcal{E}_{O}\left(  \rho\left(  \boldsymbol{\varphi}\right)
\boldsymbol{,}\mathrm{v}\right)  }{\delta\varphi_{j}^{\ast}\left(
\boldsymbol{r}\right)  }  &  =\int\frac{\delta\mathcal{E}_{O}\left(
\rho\left(  \boldsymbol{\varphi}\right)  \boldsymbol{,}\mathrm{v}\right)
}{\delta\rho\left(  \boldsymbol{r}^{\prime}\right)  }\frac{\delta\rho\left(
\boldsymbol{r}^{\prime}\right)  }{\delta\varphi_{j}^{\ast}\left(
\boldsymbol{r}\right)  }d\boldsymbol{r}^{\prime}\nonumber\\
&  =\int\frac{\delta\mathcal{E}_{O}\left(  \rho\left(  \boldsymbol{\varphi
}\right)  \boldsymbol{,}\mathrm{v}\right)  }{\delta\rho\left(  \boldsymbol{r}%
\right)  }\varphi_{j}\left(  \boldsymbol{r}^{\prime}\right)  \delta\left(
\boldsymbol{r-r}^{\prime}\right)  d\boldsymbol{r}^{\prime}=\frac
{\delta\mathcal{E}_{O}\left(  \rho\left(  \boldsymbol{\varphi}\right)
\boldsymbol{,}\mathrm{v}\right)  }{\delta\rho\left(  \boldsymbol{r}\right)
}\varphi_{j}\left(  \boldsymbol{r}\right)
\end{align}
which gives
\begin{equation}
\frac{\delta\mathcal{E}_{O}\left(  \rho\left(  \boldsymbol{\varphi}\right)
\boldsymbol{,}\mathrm{v}\right)  }{\delta\rho}\left\vert \varphi
_{j}\right\rangle =\left\vert \frac{\delta\mathcal{E}_{O}\left(  \rho\left(
\boldsymbol{\varphi}\right)  \boldsymbol{,}\mathrm{v}\right)  }{\delta\rho
}\varphi_{j}\right\rangle ,
\end{equation}
where the operator $\frac{\delta\mathcal{E}_{O}\left(  \rho\left(
\boldsymbol{\varphi}\right)  \boldsymbol{,}\mathrm{v}\right)  }{\delta\rho}$
is a multiplication operator in the coordinate representation defined by
\begin{equation}
\left[  \frac{\delta\mathcal{E}_{O}\left(  \rho\left(  \boldsymbol{\varphi
}\right)  \boldsymbol{,}\mathrm{v}\right)  }{\delta\rho}\varphi_{j}\right]
\left(  \boldsymbol{r}\right)  =\left(  \frac{\delta\mathcal{E}_{O}\left(
\rho\left(  \boldsymbol{\varphi}\right)  \boldsymbol{,}\mathrm{v}\right)
}{\delta\rho\left(  \boldsymbol{r}\right)  }\right)  \left(  \varphi
_{j}\left(  \boldsymbol{r}\right)  \right)  .
\end{equation}
Substituting this identity back into eq. $\left(  \ref{LagPar}\right)  $ one
obtains for the special case of optimized orbitals that%
\begin{equation}
\frac{\delta\mathcal{E}_{O}\left(  \rho\left(  \boldsymbol{\varphi\left(
\boldsymbol{N}\right)  }\right)  \boldsymbol{,}\mathrm{v}\right)  }{\delta
\rho}\left\vert \varphi_{j}\left(  \boldsymbol{N}\right)  \right\rangle
=\sum_{1\leq k\leq r}\left(  \varepsilon_{jk}\left(  \boldsymbol{N}\right)
-\mu\left(  N\right)  \delta_{jk}\right)  \left\vert \varphi_{k}\left(
\boldsymbol{N}\right)  \right\rangle ,\,\varepsilon_{jk}\left(  \boldsymbol{N}%
\right)  =\varepsilon_{kj}^{\ast}\left(  \boldsymbol{N}\right)  ,
\label{FockEq}%
\end{equation}
thus obtaining orbital equations corresponding to the condition eq. $\left(
\ref{KSC1}\right)  .$

The eqs $\left(  \ref{KSC2}\right)  $ and $\left(  \ref{KSC3}\right)  $
restate the constraint conditions and as $\left\Vert \left\vert \varphi
_{k}\left(  \boldsymbol{N}\right)  \right\rangle \right\Vert =N_{k}%
=0,\nu+1\leq k\leq r$ eqs $\left(  \ref{OrbLag},\ref{KSM1},\ref{KSC}%
,\ref{LagPar},\ref{FockEq}\right)  $ \emph{only} refer to the variables
$\left(  \varphi_{1},\cdots,\varphi_{\nu}\right)  $ i.e. to the occupied
orbitals. One should note that eq. $\left(  \ref{FockEq}\right)  $ is
\emph{not }an eigenvalue equation unless $\mathcal{E}_{O}\left(
\boldsymbol{\varphi}\left(  \boldsymbol{N}\right)  ,\mathrm{v}\right)  $ is
invariant to unitary mixing of $\boldsymbol{\varphi}\left(  \boldsymbol{N}%
\right)  $ that produce the same density. However to this end one can invoke
Harrimans theorem \cite{Harriman81} which shows that there is always a class
of IPS's, $\mathfrak{S}\left(  N,\rho\right)  $ (the Harriman set for
densities associated to states that contain on the average $N$ electrons) that
produce a given density, so without loss of generality we can always restrict
the domain of $\Gamma_{1}$ to $\mathfrak{S}\left(  N,\rho\right)  $ and
replace eq. $\left(  \ref{FockEq}\right)  $ by
\begin{equation}
\frac{\delta\mathcal{E}_{O}\left(  \rho\left(  \boldsymbol{\varphi\left(
\boldsymbol{N}\right)  }\right)  \boldsymbol{,}\mathrm{v}\right)  }{\delta
\rho}\left\vert \varphi_{j}\left(  \boldsymbol{N}\right)  \right\rangle
=\left(  \varepsilon_{jj}\left(  \boldsymbol{N}\right)  -\mu\left(  N\right)
\right)  \left\vert \varphi_{j}\left(  \boldsymbol{N}\right)  \right\rangle
;\quad1\leq j\leq N \label{FockHarri}%
\end{equation}
as $\mathcal{E}_{O}\left(  \boldsymbol{\varphi}\left(  \boldsymbol{\iota}%
_{N}\right)  ,\mathrm{v}\right)  $, where
\begin{equation}
\boldsymbol{\iota}_{N}=\left(  \overset{N}{\overbrace{1,\cdots,1}}%
,\overset{r-N}{\overbrace{0,\cdots,0}}\right)
\end{equation}
is inavariant to unitary mixing of $\boldsymbol{\varphi}\left(
\boldsymbol{\iota}_{N}\right)  \in\mathfrak{S}\left(  N,\rho\right)  .$

From sensitivity theory the rate of change of the orbital energy functional
with respect to changes in occupation numbers, which in this case are the
control parameters $\boldsymbol{N}$, are
\begin{equation}
\frac{\partial\mathcal{E}_{O}\left(  \boldsymbol{\varphi\left(  \boldsymbol{N}%
\right)  },\mathrm{v}\right)  }{\partial N_{j}}=\varepsilon_{jj}\left(
\boldsymbol{N}\right)  -\mu\left(  N\right)  ;1\leq j\leq\nu, \label{Janak}%
\end{equation}
which is JanaksTheorem for \emph{any }fixed set $\boldsymbol{\boldsymbol{N}}$
of occupation numbers and optimized orbitals $\boldsymbol{\varphi\left(
\boldsymbol{N}\right)  }$ normalized to these values. If we take expansion eq.
$\left(  \ref{FockEq}\right)  $ into account we can obtain these relationships
as expectation values of the operator $\frac{\delta\mathcal{E}_{O}\left(
\rho\left(  \boldsymbol{\varphi\left(  \boldsymbol{N}\right)  }\right)
\boldsymbol{,}\mathrm{v}\right)  }{\delta\rho}-\mu I$ i.e.
\begin{equation}
\frac{\partial\mathcal{E}_{O}\left(  \rho\left(  \boldsymbol{\varphi\left(
\boldsymbol{N}\right)  }\right)  ,\mathrm{v}\right)  }{\partial N_{j}%
}=\left\langle \tilde{\varphi}_{j}\left(  \boldsymbol{N}\right)  \left\vert
\left\{  \frac{\delta\mathcal{E}_{O}\left(  \rho\left(  \boldsymbol{\varphi
\left(  \boldsymbol{N}\right)  }\right)  \boldsymbol{,}\mathrm{v}\right)
}{\delta\rho}-\mu\left(  N\right)  I\right\}  \right.  \tilde{\varphi}%
_{j}\left(  \boldsymbol{N}\right)  \right\rangle =\varepsilon_{jj}\left(
\boldsymbol{N}\right)  -\mu\left(  N\right)  ,
\end{equation}
where
\begin{equation}
\left\vert \tilde{\varphi}_{j}\left(  \boldsymbol{N}\right)  \right\rangle
=\frac{\left\vert \varphi_{j}\left(  \boldsymbol{N}\right)  \right\rangle
}{\left(  \left\langle \left.  \varphi_{j}\left(  \boldsymbol{N}\right)
\right\vert \varphi_{j}\left(  \boldsymbol{N}\right)  \right\rangle \right)
^{\frac{1}{2}}}.
\end{equation}
By considering the sensitivity of the energy to changes of the average number
of electrons in the system one obtains
\begin{equation}
\frac{\partial\mathcal{E}_{O}\left(  \boldsymbol{\varphi}\left(
\boldsymbol{N}\right)  \boldsymbol{,}\mathrm{v}\right)  }{\partial N}%
=\mu\left(  N\right)
\end{equation}
and $\mu\left(  N\right)  $ can be interpreted as the chemical potential of
the interacting molecular system.

(b) In the second optimization we can restrict attention to control
parameters, $\boldsymbol{N},$ such that $\nu\left(  \boldsymbol{N}\right)
\geq N$ and $0\leq N_{j}\leq1,1\leq j\leq r$, i.e. that can be interpretted as
occupation numbers. There one finds minima of $\mathcal{E}_{O}\left(
\boldsymbol{\varphi\left(  \boldsymbol{\theta}\right)  ,}\mathrm{v}\right)
\equiv\mathcal{E}_{O}\left(  \boldsymbol{\varphi\left(  N\left(
\boldsymbol{\theta}\right)  \right)  ,}\mathrm{v}\right)  $ with respect to a
constrained variation of the normalization control parameters $\boldsymbol{N}%
=N\left(  \boldsymbol{\theta}\right)  $ parameterized by $\boldsymbol{\theta}%
$, where
\[
0\leq\theta_{j}\leq\frac{\pi}{2};\quad1\leq j\leq r
\]
for a fixed $N$ that are constrained by $\sum_{1\leq j\leq r}N_{j}=\sum_{1\leq
j\leq r}\sin\theta_{j}=N$. Stationary points $\boldsymbol{\theta}_{\#}\left(
N\right)  $ are determined by
\begin{equation}
\mathcal{L}_{\boldsymbol{\theta}}\left(  \boldsymbol{\varphi\left(
\boldsymbol{\theta}_{\#}\right)  ,}N\boldsymbol{,}\eta\left(  N\right)
\boldsymbol{,}\mathrm{v}\right)  =\max_{\boldsymbol{\eta}}\min
_{\boldsymbol{\theta}}\mathcal{L}_{\boldsymbol{\theta}}\left(
\boldsymbol{\varphi,}N\boldsymbol{,}\eta\boldsymbol{,}\mathrm{v}\right)
\end{equation}
the Lagrangian
\begin{equation}
\mathcal{L}_{\boldsymbol{\theta}}\left(  \boldsymbol{\varphi,}N\boldsymbol{,}%
\eta\boldsymbol{,}\mathrm{v}\right)  =\mathcal{E}_{O}\left(
\boldsymbol{\varphi}\left(  \boldsymbol{\theta}\right)  \boldsymbol{,}%
\mathrm{v}\right)  +\eta\left\{  \sum_{1\leq j\leq r}\sin\theta_{j}=N\right\}
\end{equation}
where
\begin{align}
\frac{\partial\mathcal{L}_{\boldsymbol{\theta}}\left(  \boldsymbol{\varphi
,}N\boldsymbol{,}\eta\boldsymbol{,}\mathrm{v}\right)  }{\partial\theta_{j}}
&  =0;1\leq j\leq r\\
\frac{\partial\mathcal{L}_{\boldsymbol{\theta}}\left(  \boldsymbol{\varphi
,}N\boldsymbol{,}\eta\boldsymbol{,}\mathrm{v}\right)  }{\partial\eta}  &
=0.\nonumber
\end{align}
The derivatives with respect to $\boldsymbol{\theta}$ are given by the chain
rule
\begin{equation}
\frac{\partial\mathcal{E}_{O}\left(  \boldsymbol{\varphi}\left(
\boldsymbol{\theta}\left(  N\right)  \right)  \boldsymbol{,}\mathrm{v}\right)
}{\partial\theta_{j}}=\sum\limits_{1\leq k\leq r}\frac{\partial\mathcal{E}%
_{O}\left(  \boldsymbol{\varphi}\left(  \boldsymbol{N}\right)  \boldsymbol{,}%
\mathrm{v}\right)  }{\partial N_{k}}\frac{\partial N_{k}}{\partial\theta_{j}%
}=\frac{\partial\mathcal{E}_{O}\left(  \boldsymbol{\varphi}\left(
\boldsymbol{N}\right)  \boldsymbol{,}\mathrm{v}\right)  }{\partial N_{j}}%
\cos\theta_{j};\quad1\leq j\leq r \label{ThetaDer}%
\end{equation}
and
\begin{equation}
\frac{\partial\mathcal{E}_{O}\left(  \boldsymbol{\varphi}\left(
\boldsymbol{N}\right)  \boldsymbol{,}\mathrm{v}\right)  }{\partial N_{j}%
}=\varepsilon_{jj}\left(  \boldsymbol{N}\left(  \boldsymbol{\theta}\right)
\right)  -\mu\left(  N\right)  .
\end{equation}
Giving
\begin{equation}
\frac{\partial\mathcal{E}_{O}\left(  \boldsymbol{\varphi}\left(
\boldsymbol{\theta}_{\#}\right)  \boldsymbol{,}\mathrm{v}\right)  }%
{\partial\theta_{j}}=\left\{  \varepsilon_{jj}\left(  \boldsymbol{N}\left(
\boldsymbol{\theta}_{\#}\right)  \right)  -\mu\left(  N\right)  \right\}
\cos\theta_{\#j}=0. \label{theta}%
\end{equation}
Thus
\begin{equation}
\varepsilon_{jj}\left(  \boldsymbol{N}\left(  \boldsymbol{\theta}_{\#}\right)
\right)  -\mu\left(  N\right)  =0\text{ and/or }\left\{  \cos\theta
_{\#j}=0\Leftrightarrow\theta_{\#j}=\frac{\pi}{2}\right\}  ;\quad1\leq j\leq
r. \label{Con1}%
\end{equation}
If $\ \theta_{j}=\frac{\pi}{2}$ ($N_{j}\left(  \boldsymbol{\theta}%
_{\#}\right)  =1)$ then $\varepsilon_{jj}$ can have any value (even
$\mu\left(  N\right)  $), while if $\theta_{\#j}\neq\frac{\pi}{2}$ (which
corresponds to non integer occupation) then $\varepsilon_{jj}=$ $\mu\left(
N\right)  $ i.e. the Lagrange parameters $\left\{  \varepsilon_{jj}\right\}  $
are degenerate. If $\varepsilon_{jj}=\mu\left(  N\right)  $ then $\theta
_{\#j}$ does not have to have the value $\frac{\pi}{2}$ though it can. Since
the inception of the KS approach to DFT\ there has been an extensive
discussion of the physical meaning of the eigenvalues of the KS operator and
their relationship to the Fermi energy and chemical potential.

The ground state energy $E\left(  \rho_{GS}\boldsymbol{,}\mathrm{v}\right)
=\mathcal{E}_{O}\left(  \boldsymbol{\varphi}\left(  \boldsymbol{\theta}%
_{\#}\right)  ,\mathrm{v}\right)  $ and the ground state density can then be
expanded as
\begin{equation}
\rho_{GS}\left(  \boldsymbol{r}\right)  =\sum_{1\leq j\leq r}\left\vert
\left[  \varphi_{j}\left(  \boldsymbol{\theta}_{\#}\right)  \right]  \left(
\boldsymbol{r}\right)  \right\vert ^{2}. \label{GSen}%
\end{equation}
The expansion eq. $\left(  \ref{GSen}\right)  $ is very \emph{non} unique in
that any set $\boldsymbol{\varphi}\left(  \boldsymbol{\theta}_{\#}^{\prime
}\right)  \in\mathfrak{S}\left(  \rho_{GS}\right)  $, where
$\boldsymbol{\theta}_{\#}^{\prime}$ satisfies eq. $\left(  \ref{theta}\right)
$ produces $\rho_{GS}$ and the same value of the energy functional i.e.
$\mathcal{E}_{O}\left(  \boldsymbol{\varphi}\left(  \boldsymbol{\theta}%
_{\#}^{\prime}\right)  ,\mathrm{v}\right)  =\mathcal{E}_{O}\left(
\boldsymbol{\varphi}\left(  \boldsymbol{\theta}_{\#}\right)  ,\mathrm{v}%
\right)  $.

\subsection{The Kohn-Sham Density Functional\label{KSDF}}

As pointed out in the last subsection the minima of $\mathcal{E}_{O}\left(
\boldsymbol{\varphi}\left(  \boldsymbol{\theta}_{\#}\right)  ,\mathrm{v}%
\right)  ,$
\begin{equation}
\mathcal{E}_{O}\left(  \boldsymbol{\varphi}\left(  \boldsymbol{\theta}%
_{\#}\right)  \boldsymbol{,}\mathrm{v}\right)  =\min_{\boldsymbol{\varphi\in
}\mathfrak{S}_{1}}\mathcal{E}_{O}\left(  \boldsymbol{\varphi,}\mathrm{v}%
\right)
\end{equation}
are highly degenerate e.g.
\begin{equation}
\mathcal{E}_{O}\left(  \boldsymbol{\varphi}\left(  \boldsymbol{\theta}%
_{\#}\right)  \boldsymbol{,}\mathrm{v}\right)  =const.\,\forall
\boldsymbol{\varphi\in}\mathfrak{S}\left(  \rho_{GS}\right)  .
\end{equation}
This degeneracy is due to the singularity of the map $\Gamma_{1}$ which can be
removed by constraining $\boldsymbol{\varphi}$ to be a specific representative
of the equivalence classes $\mathfrak{S}\left(  \rho\right)  .$ One way to
obtain such a representative is to consider the orbitals formed by the NOsand
occupation numbers of the FORDO of the $N$-electron ground state $D_{0}\left(
\rho\right)  =A_{0}\left(  \rho\right)  ^{2}$ of a model independent particle
$N$-electron system$,$ that has a fixed density $\rho.$ This corresponds to
finding the minimizers of the universal kinetic energy functional $K_{U}$
\begin{equation}
K_{U}\left(  A\right)  =\operatorname{Tr}\left\{  KA^{2}\right\}  ,
\end{equation}
subject to the constraint
\begin{equation}
\Xi_{2}\left(  A\right)  =\left(  A\left\vert \widehat{\Phi^{\dagger}\Phi
}A\right.  \right)  =\rho,\,A\in\mathcal{B}_{2sa}\left(  \mathcal{H}%
^{N}\right)  .
\end{equation}
The constrained minima $K_{U}\left(  D_{0}\left(  \rho\right)  \right)
=\operatorname{Tr}\left\{  KD_{0}\left(  \rho\right)  \right\}
=\operatorname{Tr}\left\{  KA_{0}\left(  \rho\right)  ^{2}\right\}  $ is given
as an unconstrained saddle point $\left(  A_{0}\left(  \rho\right)
,\lambda_{0}\left(  \rho\right)  \right)  ,$
\begin{equation}
\mathcal{L}_{K}\left(  A_{0}\left(  \rho\right)  ,\lambda_{0}\left(
\rho\right)  ,\rho\right)  =\max_{\lambda}\min_{A}\mathcal{L}_{K}\left(
A,\lambda,\rho\right)
\end{equation}
of the Lagrangian
\begin{equation}
\mathcal{L}_{K}\left(  A,\lambda,\rho\right)  =K_{U}\left(  A\right)
-\Omega_{A}\left(  \lambda\right)  +\left\langle \left.  \lambda\right\vert
\rho\right\rangle =\left(  A\left\vert \left\{  \left\{  \widehat{K}%
-\widehat{\Omega\left(  \lambda\right)  }\right\}  +\left\langle \left.
\lambda\right\vert \rho\right\rangle \hat{I}\right\}  \right.  A\right)  .
\end{equation}
If a bound state minimum exists for some local one body potential $\lambda
_{0}\left(  \rho\right)  $ (which is also the Lagrange parameter function) one
has
\begin{equation}
\left\{  \widehat{K}-\widehat{\Omega\left(  \lambda_{0}\left(  \rho\right)
\right)  }\right\}  \left\vert A_{0}\left(  \rho\right)  \right)  =0,
\end{equation}
where $A_{0}\left(  \rho\right)  \in\mathcal{B}_{2sa}\left(  \mathcal{H}%
^{N}\right)  $. This gives a stationary state $D_{0}\left(  \rho\right)
=A_{0}\left(  \rho\right)  ^{2}$ of the independent particle system with
Hamiltonian $\left\{  K-\Omega\left(  \lambda_{0}\left(  \rho\right)  \right)
\right\}  $ i.e.
\begin{equation}
\left[  \left\{  K-\Omega\left(  \lambda_{0}\left(  \rho\right)  \right)
\right\}  ,D_{0}\left(  \rho\right)  \right]  =0,
\end{equation}
which in turn defines a FORDO $D_{0}^{1}\left(  \rho\right)  $ given by (see
Appendix \ref{Denmat})
\begin{equation}
D_{0}^{1}\left(  \rho\right)  =NL_{N}^{1}\left(  D_{0}\left(  \rho\right)
\right)  .
\end{equation}
As $K-\Omega\left(  \lambda_{0}\left(  \rho\right)  \right)  $ is an
one-electron operator the state $D_{0}\left(  \rho\right)  $ must be a mixture
of IPSsor a pure IPS and thus determined by the eigenvalues $\left\{
N_{0j}\left(  \rho\right)  \right\}  $ and eigenvectors $\left\{  \left\vert
\widetilde{\varphi_{0j}\left(  \rho\right)  }\right\rangle \right\}  $ of
$D_{0}^{1}\left(  \rho\right)  $. If $D_{0}\left(  \rho\right)  $ is a pure
IPS\ then the occupation numbers $\left\{  N_{0j}\left(  \rho\right)
\right\}  $ are equal to either $1$ or $0$ (i.e. $N$ ones and $r-N$ zeroes),
while if it is a ensemble state than there will be some fractional values and
the number of zero values will be less than $r-N$. Expanding $D_{0}^{1}\left(
\rho\right)  $ in terms of normalized NOs$\left\{  \left.  \left\vert
\widetilde{\varphi_{0j}}\right\rangle \right\vert 1\leq j\leq r\right\}  $
gives
\begin{equation}
D_{0}^{1}\left(  \rho\right)  =\sum_{1\leq j\leq r}N_{0j}\left(  \rho\right)
\left\vert \widetilde{\varphi_{0j}\left(  \rho\right)  }\right\rangle
\left\langle \widetilde{\varphi_{0j}\left(  \rho\right)  }\right\vert ,
\end{equation}
(where in contrast to eq. $\left(  \ref{Gamma}\right)  $ $r\geq N$ ), which
after renormalizing the NOsto the occupation numbers i.e.
\begin{equation}
\left\vert \varphi_{0j}\left(  \rho\right)  \right\rangle =N_{0j}\left(
\rho\right)  ^{\frac{1}{2}}\left\vert \widetilde{\varphi_{0j}\left(
\rho\right)  }\right\rangle ,1\leq j\leq r
\end{equation}
results in
\begin{equation}
\rho\left(  \boldsymbol{r}\right)  =\sum_{1\leq j\leq r}\left\vert \left[
\varphi_{0j}\left(  \rho\right)  \right]  \left(  \boldsymbol{r}\right)
\right\vert ^{2}. \label{gammaKS}%
\end{equation}
If the constrained minimum of $K_{U}$ is non degenerate then there is a unique
correspondence, up to unitary mixing of degenerate eigenvectors of $D_{0}%
^{1}\left(  \rho\right)  ,$ between $\rho$ and $\left\{  \boldsymbol{\varphi
}_{0}\right\}  ,$ thus modulo this unitary mixing one can define an inverse
functional $\Gamma_{0}^{-1}$ as
\begin{equation}
\boldsymbol{\varphi}_{0}=\Gamma_{0}^{-1}\left(  \rho\right)  .
\end{equation}
If the constrained minimum of $K_{U}$ is degenerate, e.g. corresponding to
FORDOs$\left\{  \left.  D_{0\kappa}^{1}\left(  \rho\right)  \right\vert
1\leq\kappa\leq m\right\}  $ of the stationary states $\left\{  \left.
D_{0\kappa}\left(  \rho\right)  \right\vert 1\leq\kappa\leq m\right\}  $ of
$K-\Omega\left(  \lambda_{0}\left(  \rho\right)  \right)  $ then as
\begin{equation}
\left[  \left\{  K-\Omega\left(  \lambda_{0}\left(  \rho\right)  \right)
\right\}  ,D_{0\kappa}\left(  \rho\right)  \right]  =\left[  \left\{
K-\Omega\left(  \lambda_{0}\left(  \rho\right)  \right)  \right\}
,D_{0\kappa}^{1}\left(  \rho\right)  \right]  =0,
\end{equation}
the space of degenerate states corresponds to the same set of orbitals
$\left\{  \left\vert \boldsymbol{\varphi}_{0}\left(  \rho\right)
\right\rangle \right\}  $ but normalized to different occupation numbers
$\left\{  \boldsymbol{N}_{0}\left(  \rho\right)  \right\}  ,$ where
$\boldsymbol{N}_{0}\left(  \rho\right)  $ are convex combinations of $\left\{
\left.  \boldsymbol{N}_{0\kappa}\left(  \rho\right)  \right\vert 1\leq
\kappa\leq m\right\}  $. Hence if degeneracies occur these are reflected in
multiple optimal values for $\boldsymbol{\boldsymbol{N}}_{0}$ for fixed
$\left\{  \boldsymbol{\varphi}_{0}\left(  \rho\right)  \right\}  $, some of
the $\left\{  \left.  N_{0j}\right\vert 1\leq j\leq r\right\}  $ can have non
integer values even for degenerate pure states$\boldsymbol{.}$

The minimal values of the universal kinetic functional $K_{U}\left(
D_{0}\left(  \rho\right)  \right)  $ produce a kinetic energy functional
$\mathcal{F}_{T}$ defined on the representatives $\left\{  \boldsymbol{\varphi
}_{0}\right\}  $ of the equivalence classes $\mathfrak{S}\left(  \rho\right)
$ whose values are given by
\begin{equation}
\mathcal{F}_{T}\left(  \boldsymbol{\varphi}_{0}\right)  =K_{U}\left(
D_{0}\left(  \Gamma_{1}\left(  \boldsymbol{\varphi}_{0}\right)  \right)
\right)  . \label{KEFunc}%
\end{equation}
The value of $\mathcal{F}_{T}$ is constant on the set of orbitals formed from
the unitary mixing and the convex combinations of degenerate states just
described. This allows one to define the Kohn-Sham functional $\mathcal{F}%
_{KS}\left(  \boldsymbol{\varphi}_{0}\right)  $ as a specific decomposition of
the HK energy functional $\mathcal{F}_{HK}\left(  \rho\right)  $ in terms of
orbitals as
\begin{equation}
\mathcal{F}_{KS}\left(  \boldsymbol{\varphi}_{0}\right)  =\mathcal{F}%
_{T}\left(  \boldsymbol{\varphi}_{0}\right)  +\mathcal{F}_{C}\left(
\boldsymbol{\varphi}_{0}\right)  +\mathcal{F}_{KSXC}\left(
\boldsymbol{\varphi}_{0}\right)  ,
\end{equation}
where
\begin{equation}
\mathcal{F}_{KS}\left(  \boldsymbol{\varphi}_{0}\right)  =\mathcal{F}%
_{HK}\left(  \Gamma_{1}\left(  \boldsymbol{\varphi}_{0}\right)  \right)  .
\end{equation}


The relationships described in subsection \ref{GSOFunc} derived for stationary
points $\boldsymbol{\varphi}\left(  \boldsymbol{N}\right)  $ are valid for any
points contained in the equivalence classes that these points belong to, in
particular $\boldsymbol{\varphi}_{0}\left(  \boldsymbol{\boldsymbol{N}}%
_{0}\right)  $. Eq. $\left(  \ref{FockEq}\right)  $, is now an equation for
determining the KS orbitals $\boldsymbol{\varphi}_{0}\left(  \boldsymbol{N}%
\right)  $ and the operator introduced there can be identified as the KS
operator. If the state $D_{0}\left(  \rho\right)  $ is non degenerate and pure
(therefore has only occupation numbers equal to one or zero) then one can
choose a set of orbitals that diagonalize this operator as the energy
$\mathcal{E}_{O}\left(  \boldsymbol{\varphi}_{0}\left(  \boldsymbol{N}%
_{0}\right)  ,\mathrm{v}\right)  $ is invariant to unitary mixing of occupied
orbitals, which is the standard way of producing a canonical set of occupied orbitals.

The reference $N$-electron IPS, $D_{0}\left(  \boldsymbol{\boldsymbol{N}}%
_{0}\right)  $ formed from the NOs$\boldsymbol{\varphi}_{0}$ has two energies
associated with it. One, $E_{0}\left(  \boldsymbol{\varphi}_{0}\left(
\boldsymbol{N}_{0}\right)  \right)  $, is the ground state energy of the
Hamiltonian%
\begin{equation}
H_{0}=K+\mathrm{v},
\end{equation}
which can be expressed as
\begin{equation}
E_{0}\left(  \boldsymbol{\varphi}_{0}\left(  \boldsymbol{N}_{0}\right)
\right)  =\sum_{1\leq j\leq r}\left\langle \tilde{\varphi}_{0j}\left\vert
\left(  K+\mathrm{v}\right)  \right.  \tilde{\varphi}_{0j}\right\rangle
N_{0j}.
\end{equation}
The other is the energy of the Hamiltonian
\begin{equation}
H=H_{0}+V_{2}%
\end{equation}
wrt the IPS $D_{0}\left(  \boldsymbol{N}_{0}\right)  $ given by
\begin{equation}
E_{IP}\left(  D_{0}\left(  \boldsymbol{N}_{0}\right)  \right)
=\operatorname{Tr}\left\{  HD_{0}\left(  \boldsymbol{N}_{0}\right)  \right\}
.
\end{equation}
It should be noted that this value of the energy functional $E_{IP}$ is
\emph{not }optimized in any way and is \emph{not }related in any simple
fashion to the energies $\mathcal{E}_{O}\left(  \varphi\boldsymbol{\left(
\boldsymbol{N}_{0}\right)  },\mathrm{v}\right)  .$

The preceding relationships are equally for orbitals $\boldsymbol{\varphi}%
_{0}\left(  \boldsymbol{N}_{0}\left(  \boldsymbol{\theta}_{\#}\left(
N\right)  \right)  \right)  $ that produce the ground state density $\rho
_{GS}$ and energy $\mathcal{E}_{O}\left(  \boldsymbol{\varphi}_{0}\left(
\boldsymbol{N}\left(  \boldsymbol{\theta}_{\#}\left(  N\right)  \right)
\right)  ,\mathrm{v}\right)  $ and the classical Janaks Theorem is a special
case of eq. $\left(  \ref{Janak}\right)  $
\begin{equation}
\frac{\partial\mathcal{E}_{O}\left(  \boldsymbol{\varphi}_{0}\left(
\boldsymbol{N}_{0}\left(  \boldsymbol{\theta}_{\#}\left(  N\right)  \right)
\right)  ,\mathrm{v}\right)  }{\partial N_{j}}\equiv\frac{\partial
\mathcal{E}_{KS}\left(  \boldsymbol{\varphi}_{0}\left(  \boldsymbol{N}%
_{0}\left(  \boldsymbol{\theta}_{\#}\left(  N\right)  \right)  \right)
,\mathrm{v}\right)  }{\partial N_{j}}=\varepsilon_{jj}\left(
\boldsymbol{\varphi}_{0}\left(  \boldsymbol{N}_{0}\left(  \boldsymbol{\theta
}_{\#}\left(  N\right)  \right)  \right)  \right)  , \label{ClassJanak}%
\end{equation}
where we have constrained the orbitals and occupation numbers to be the ones
that minimize the universal kinetic energy, $K_{U},$ thus producing a model
IPS. One must, however, note that the Lagrange parameters $\left\{
\varepsilon_{jj}\left(  \boldsymbol{\varphi}_{0}\left(  \boldsymbol{N}%
_{0}\left(  \boldsymbol{\theta}_{\#}\left(  N\right)  \right)  \right)
\right)  \right\}  $ in this equation \emph{are not }associated with the
minimization of the universal kinetic energy.

The KS energy functional is defined in terms of orbitals and occupation
numbers and is implicitly defined by the preceding constrained optimization
procedures, that factor the functional through the density. If one did not
consider this factorization and directly optimized the orbitals and occupation
numbers one would determine the same ground state, but not need to introduce a
model IPS\ and orbitals $\boldsymbol{\varphi}_{0}\left(  \boldsymbol{N}%
_{0}\right)  $. In the next section we compare and contrast the relationships
obtained for these functionals with ones obtained for the explicitly defined
AGP functional. These functionals are defined on the space of correlated
states formed by the non linear manifold of AGP\ states, where the energy has
been expressed as an explicit functional of GSOsand parameters that determine
occupation numbers \cite{OrtizWeiner2006} The minimization of this explicit
functional can then be compared to the minimization of the implicit
KS\ functional restricted to the AGP\ manifold.

\section{AGP Energy Functionals \label{AGP}}

We briefly review some of the results found in \cite{OrtizWeiner2006}
describing AGP states and their associated energy functionals. In
\cite{OrtizWeiner2006} orbital and occupation number optimization equations
were derived and it was displayed how the $N$-electron total energy could be
expressed in terms of ionization potentials, kinetic energies and eigenvalues
of a generalized Fock operator associated with GSO's. In contrast to the KS
approach, the occupation numbers were not constrained to refer to IPS's, but
allowed to vary over the much larger domain of states with doubly degenerate
state occupation numbers. It has been argued while this set does not contain
all possible states it contains a large percentage of states important for
molecular systems \cite{Kramers Deg}. The density matrix energy functional in
DMFT allows unrestricted variation of occupation numbers (and thus all
states), but unlike the AGP functional (but in common with the HK\ and KS
density functionals) it is invariant to the phases of the GSOsi.e. the group
$U\left(  1\right)  .$

\subsection{Canonical expansion of an AGP state}

The $2M$-electron AGP state, $\left\vert g^{M}\right\rangle $, is the $M^{th}$
antisymmetric power of the $2$-electron state described by a normalized
geminal
\begin{equation}
\left\vert g\right\rangle =\sum_{1\leq j\leq s}\left\vert \varphi_{j}%
\varphi_{j+s}\right\rangle =\sum_{1\leq j\leq s}n_{j}^{\frac{1}{2}}\left\vert
\tilde{\varphi}_{j}\tilde{\varphi}_{j+s}\right\rangle , \label{geminal}%
\end{equation}
given by
\begin{equation}
\left\vert g^{M}\right\rangle =\left(  S_{M}\right)  ^{-\frac{1}{2}}%
\sum_{1\leq j_{1}<\cdots<j_{M}\leq s}n_{j_{1}}^{\frac{1}{2}}\cdots n_{j_{s}%
}^{\frac{1}{2}}\left\vert \tilde{\varphi}_{j_{1}}\tilde{\varphi}_{j_{1}%
+s}\cdots\tilde{\varphi}_{j_{M}}\tilde{\varphi}_{j_{M}+s}\right\rangle
=\left(  S_{M}\right)  ^{-\frac{1}{2}}\sum_{1\leq j_{1}<\cdots<j_{M}\leq
s}\left\vert \varphi_{j_{1}}\varphi_{j_{1}+s}\cdots\varphi_{j_{M}}%
\varphi_{j_{M}+s}\right\rangle , \label{agp}%
\end{equation}
wherer $\left\{  \left.  n_{j}^{\frac{1}{2}}\right\vert 1\leq j\leq s\right\}
$ are the positive square roots of the occupation numbers%
\begin{equation}
\Omega=\left\{  n_{j}\left\vert 0\leq n_{j}\leq1;\sum_{1\leq j\leq s}%
n_{j}=1;1\leq j\leq s\right.  \right\}
\end{equation}
of the FORDO of the two electron AGP\ state density operator, $\left\{
\left.  \left\vert \tilde{\varphi}_{j}\right\rangle \right\vert 1\leq
j\leq2s\right\}  $ are the \emph{orthonormal} Canonical General Spin Orbitals
(CGSO's) of $\left\vert g\right\rangle $ and $\left\{  \left.  \left\vert
\varphi_{j}\right\rangle \right\vert 1\leq j\leq2s\right\}  $ are the
\emph{orthogonal} CGSOsnormalized to $\boldsymbol{n}$ i.e.
\begin{equation}
\left\vert \varphi_{j}\right\rangle =n_{j}^{\frac{1}{2}}\left\vert
\tilde{\varphi}_{j}\right\rangle ,1\leq j\leq2s.
\end{equation}
The normalization factor, $\left(  S_{M}\right)  ^{-\frac{1}{2}}$, is given in
terms of the elementary symmetric function
\begin{equation}
S_{M}=\sum_{1\leq j_{1}<\cdots<j_{M}\leq s}n_{j_{1}}\cdots n_{j_{M}}.
\end{equation}
It is evident from eq.\ $\left(  \ref{geminal}\right)  $ that the geminal and
geminal power state are at least invariant to the group $SU\left(  2\right)
,$ that is formed by $2\times2$ \emph{special} unitary transformations of the
paired CGSOs$\left\{  \left.  \left\vert \tilde{\varphi}_{j}\right\rangle
,\left\vert \tilde{\varphi}_{j+s}\right\rangle \right\vert 1\leq j\leq
s\right\}  ,$ while if degeneracies exist amongst the $\left\{  \left.
n_{j}\right\vert 1\leq j\leq s\right\}  $ this invariance group is larger
\cite{AGPCoherent}. The canonical CGSOsare unique up to transformations
belonging this invariance group.

\subsection{AGP Energy Functionals\label{AGPSect}}

The energy of the AGP state can be expressed as a functional $\mathcal{E}%
_{AGP}$ of $\boldsymbol{\varphi}\left(  \boldsymbol{n}\right)  $ by
\begin{align}
\mathcal{E}_{AGP}\left(  \boldsymbol{\varphi}\left(  \boldsymbol{n}\right)
\right)   &  =\frac{\mathcal{H}\left(  \boldsymbol{\varphi}\left(
\boldsymbol{n}\right)  \right)  }{S_{M}\left(  \boldsymbol{n}\right)
}\nonumber\\
&  =\frac{1}{S_{M}\left(  \boldsymbol{n}\right)  }\left\{  \sum_{1\leq j\leq
s}n_{j}S_{M-1}\left(  \jmath\right)  \left\{  \left\langle \tilde{\varphi}%
_{j}\left\vert h_{1}\right.  \tilde{\varphi}_{j}\right\rangle +\left\langle
\tilde{\varphi}_{j+s}\left\vert h\right.  \tilde{\varphi}_{j+s}\right\rangle
+\left\langle \tilde{\varphi}_{j}\tilde{\varphi}_{j+s}||\tilde{\varphi}%
_{j}\tilde{\varphi}_{j+s}\right\rangle \right\}  \right. \nonumber\\
&  +\sum_{1\leq j,k\leq s}n_{j}^{\frac{1}{2}}n_{k}^{\frac{1}{2}}S_{M-1}\left(
\jmath k\right)  \left\langle \tilde{\varphi}_{j}\tilde{\varphi}_{j+s}%
||\tilde{\varphi}_{k}\tilde{\varphi}_{k+s}\right\rangle \left.  +\sum_{1\leq
j<k\leq s}n_{i}n_{j}S_{M-2}\left(  \jmath k\right)  \left\langle
\tilde{\varphi}_{j}\tilde{\varphi}_{k}|||\tilde{\varphi}_{j}\tilde{\varphi
}_{k}\right\rangle \right\}  . \label{energy}%
\end{align}


\subsection{AGP Lagrangians\label{Lagran}}

Using the same philosophy for optimization as discussed for orbital
functionals in the preceding sections one can minimize the energy of the
$2M$-electron state over the manifold of AGP\ states using a nested two step process:

\begin{enumerate}
\item Keeping the normalization $\boldsymbol{n}$ of the
GSOs$\boldsymbol{\varphi}$ fixed one finds the constrained minimum
$\mathcal{E}_{AGP}\left(  \boldsymbol{\varphi}\left(  \boldsymbol{n}\right)
\right)  $ of the functional $\mathcal{E}_{AGP}\left(  \boldsymbol{\varphi
}\right)  $ for the specified values of the control parameters $\boldsymbol{n}%
$ i.e.
\begin{equation}
\mathcal{E}_{AGP}\left(  \boldsymbol{\varphi}\left(  \boldsymbol{n}\right)
\right)  =\min_{\boldsymbol{\varphi\in}\Lambda\left(  \boldsymbol{n}\right)
}\mathcal{E}_{AGP}\left(  \boldsymbol{\varphi}\right)  ,
\end{equation}
where
\[
\Lambda\left(  \boldsymbol{n}\right)  =\left\{  \boldsymbol{\varphi}\left\vert
\left\langle \left.  \varphi_{j}\right\vert \varphi_{k}\right\rangle
=n_{j}\delta_{jk};\sum_{1\leq j\leq s}n_{j}=1\right.  \right\}  .
\]
This is accomplished by finding the unconstrained saddle points $\mathcal{L}%
\left(  \boldsymbol{\varphi\left(  \boldsymbol{n}\right)  ,\varepsilon\left(
\boldsymbol{n}\right)  },\eta\right)  $ of the Lagrangian
\begin{equation}
\mathcal{L}\left(  \boldsymbol{\varphi,\varepsilon}\right)  =\mathcal{E}%
_{AGP}\left(  \boldsymbol{\varphi}\right)  -\sum_{1\leq j,k\leq2s}%
\varepsilon_{kj}\left(  \left\langle \left.  \varphi_{j}\right\vert
\varphi_{k}\right\rangle -n_{j}\delta_{jk}\right)  +\eta\left\{  \sum_{\,1\leq
j\leq2s}\left\langle \left.  \varphi_{j}\right\vert \varphi_{j}\right\rangle
-2\right\}  ,
\end{equation}
i.e.
\begin{equation}
\mathcal{L}\left(  \boldsymbol{\varphi\left(  \boldsymbol{n}\right)
,\varepsilon\left(  \boldsymbol{n}\right)  },\eta\right)  =\max
_{\boldsymbol{\varepsilon},\eta.}\min_{\boldsymbol{\varphi}}\mathcal{L}\left(
\boldsymbol{\varphi,\varepsilon},\eta\right)
\end{equation}
The corresponding stationarity conditions are
\begin{subequations}
\label{AGPstat}%
\begin{align}
\frac{\delta\mathcal{L}\left(  \boldsymbol{\varphi\left(  \boldsymbol{n}%
\right)  ,\varepsilon\left(  \boldsymbol{n}\right)  },\eta\right)  }%
{\delta\boldsymbol{\varphi}}  &  =\boldsymbol{0}\label{agpstat1}\\
\frac{\partial\mathcal{L}\left(  \boldsymbol{\varphi\left(  \boldsymbol{n}%
\right)  ,\varepsilon\left(  \boldsymbol{n}\right)  },\eta\right)  }%
{\partial\boldsymbol{\varepsilon}}  &  =\boldsymbol{0}\label{agpstat2}\\
\frac{\partial\mathcal{L}\left(  \boldsymbol{\varphi\left(  \boldsymbol{n}%
\right)  ,\varepsilon\left(  \boldsymbol{n}\right)  },\eta\right)  }%
{\partial\eta}  &  =\boldsymbol{0} \label{agpstat3}%
\end{align}
and eq. $\left(  \ref{agpstat1}\right)  $ gives
\end{subequations}
\begin{equation}
\frac{\delta\mathcal{E}_{AGP}\left(  \boldsymbol{\varphi}\left(
\boldsymbol{n}\right)  \right)  }{\delta\varphi_{j}}=\sum_{1\leq k\leq
2s}\varepsilon\boldsymbol{\left(  \boldsymbol{n}\right)  }_{jk}\left\vert
\varphi_{k}\boldsymbol{\left(  \boldsymbol{n}\right)  }\right\rangle ,
\end{equation}
which can be expressed \cite{OrtizWeiner2006} in terms of a generalized Fock
operator as
\begin{equation}
F\left(  \boldsymbol{\varphi}\left(  \boldsymbol{n}\right)  \right)
\left\vert \varphi_{j}\left(  \boldsymbol{n}\right)  \right\rangle
=\sum_{1\leq k\leq2s}\varepsilon\boldsymbol{\left(  \boldsymbol{n}\right)
}_{jk}\left\vert \varphi_{k}\boldsymbol{\left(  \boldsymbol{n}\right)
}\right\rangle .
\end{equation}
Sensitivity theory gives%
\[
\frac{\partial\mathcal{E}_{AGP}\left(  \boldsymbol{\varphi}\left(
\boldsymbol{n}\right)  \right)  }{\partial n_{j}}=\varepsilon_{jj}\left(
\boldsymbol{n}\right)  .
\]


\item Varying $\mathcal{E}_{AGP}\left(  \boldsymbol{\varphi}\left(
\boldsymbol{n}\right)  \right)  $ over the normalization control parameters
$\boldsymbol{n}\in\Omega$ to obtain the minimum $\mathcal{E}_{AGP}\left(
\boldsymbol{\varphi}\left(  \boldsymbol{n}_{\#}\right)  \right)  $ i.e.
\begin{equation}
\mathcal{E}_{AGP}\left(  \boldsymbol{\varphi}\left(  \boldsymbol{n}%
_{\#}\right)  \right)  =\min_{\boldsymbol{n}\in\Omega}\mathcal{E}_{AGP}\left(
\boldsymbol{\varphi}\left(  \boldsymbol{n}\right)  \right)  .
\label{agp_control}%
\end{equation}
The optimization eq. $\left(  \ref{agp_control}\right)  $ can be made
unconstrained by expressing $\boldsymbol{n}$ in terms of the variable
$\boldsymbol{\theta}$ as
\begin{equation}
n_{j}\left(  \boldsymbol{\theta}\right)  =\frac{\cos^{2}\theta_{j}}%
{\sum\limits_{1\leq j\leq s}\cos^{2}\theta_{j}}\equiv\frac{\cos^{2}\theta_{j}%
}{\mathcal{N}_{n}\left(  \boldsymbol{\theta}\right)  };0\leq\theta_{j}%
\leq\frac{\pi}{2} \label{ntheta}%
\end{equation}
$\lceil$%
\begin{equation}
\sum\limits_{1\leq j\leq s}\frac{\partial\mathcal{E}_{AGP}\left(
\boldsymbol{\varphi}\left(  \boldsymbol{n}\right)  \right)  }{\partial n_{j}%
}\frac{\partial n_{j}}{\partial\theta_{k}}=0;1\leq k\leq s.
\end{equation}
Using eq.$\left(  \ref{ntheta}\right)  $ and observing that
\begin{equation}
\frac{\partial n_{j}}{\partial\theta_{k}}=\frac{\cos^{2}\theta_{j}\sin
2\theta_{k}}{\mathcal{N}\left(  \boldsymbol{\theta}\right)  ^{2}}-\frac
{\sin2\theta_{k}\delta_{jk}}{\mathcal{N}\left(  \boldsymbol{\theta}\right)
}=\frac{\sin2\theta_{k}}{\mathcal{N}\left(  \boldsymbol{\theta}\right)
}\left\{  \frac{\cos^{2}\theta_{j}}{\mathcal{N}\left(  \boldsymbol{\theta
}\right)  }-\delta_{jk}\right\}
\end{equation}
and
\begin{equation}
\frac{\partial\mathcal{E}_{AGP}\left(  \boldsymbol{\varphi}\left(
\boldsymbol{n}\left(  \boldsymbol{\theta}\right)  \right)  \right)  }{\partial
n_{j}}=\varepsilon_{jj}\left(  \boldsymbol{n}\left(  \boldsymbol{\theta
}\right)  \right)
\end{equation}
this can be developed to give
\begin{align}
&  \frac{\partial\mathcal{E}_{AGP}\left(  \boldsymbol{\varphi}\left(
\boldsymbol{\theta}\right)  \right)  }{\partial\theta_{k}}=\sum\limits_{1\leq
j\leq s}\varepsilon_{jj}\left(  \boldsymbol{\theta}\right)  \frac{\sin
2\theta_{k}}{\mathcal{N}\left(  \boldsymbol{\theta}\right)  }\left\{
\frac{\cos^{2}\theta_{j}}{\mathcal{N}\left(  \boldsymbol{\theta}\right)
}-\delta_{jk}\right\}  =0\nonumber\\
&  \Rightarrow\sum\limits_{1\leq j\leq s}\varepsilon_{jj}\left(
\boldsymbol{\theta}\right)  \frac{\cos^{2}\theta_{j}\sin2\theta_{k}%
}{\mathcal{N}\left(  \boldsymbol{\theta}\right)  ^{2}}=\varepsilon_{kk}\left(
\boldsymbol{\theta}\right)  \sin2\theta_{k}%
\end{align}
$\rfloor$
\end{enumerate}

We are interested in the dependence on the occupation numbers of the
$2M$-electron state which we can obtain by using the chain rule
\begin{align}
&  \frac{\partial\mathcal{E}_{AGP}\left(  \boldsymbol{\varphi}\left(
\boldsymbol{n}\right)  \right)  }{\partial N_{j}}=\sum_{1\leq k\leq2s}%
\frac{\partial\mathcal{E}_{AGP}\left(  \boldsymbol{\varphi}\left(
\boldsymbol{n}\right)  \right)  }{\partial n_{k}}\frac{\partial n_{k}%
}{\partial N_{j}}\nonumber\\
&  \Rightarrow\frac{\partial\mathcal{E}_{AGP}\left(  \boldsymbol{\varphi
}\left(  \boldsymbol{n}\right)  \right)  }{\partial N_{j}}=\sum_{1\leq
k\leq2s}\frac{\partial\mathcal{E}_{AGP}\left(  \boldsymbol{\varphi}\left(
\boldsymbol{n}\right)  \right)  }{\partial n_{k}}\frac{\partial n_{k}%
}{\partial N_{j}}=\sum_{1\leq k\leq2s}\frac{\partial\mathcal{E}_{AGP}\left(
\boldsymbol{\varphi}\left(  \boldsymbol{n}\right)  \right)  }{\partial n_{k}%
}\left(  \frac{\partial N_{k}}{\partial n_{j}}\right)  ^{-1}%
\end{align}
i.e.%
\begin{equation}
\frac{\partial\mathcal{E}_{AGP}\left(  \boldsymbol{\varphi}\left(
\boldsymbol{n}\right)  \right)  }{\partial\boldsymbol{N}}=\frac{\partial
\mathcal{E}_{AGP}\left(  \boldsymbol{\varphi}\left(  \boldsymbol{n}\right)
\right)  }{\partial\boldsymbol{n}}\left(  \frac{\partial\boldsymbol{N}%
}{\partial\boldsymbol{n}}\right)  ^{-1}.
\end{equation}
The Jacobian $\frac{\partial\boldsymbol{N}}{\partial\boldsymbol{n}}$ is
nonsingular as the occupation numbers of the AGP\ state are explicitly given
in a $1-1$ fashion by $\left[  {}\right]  $
\begin{equation}
\boldsymbol{N}=\mathcal{N}\left(  \boldsymbol{n}\right)  ,
\end{equation}
where
\begin{equation}
N_{j}=\mathcal{N}_{j}\left(  \boldsymbol{n}\right)  =n_{j}\frac{S_{M-1}\left(
\jmath\right)  }{S_{M}}=n_{j}\frac{\partial\ln S_{M}}{\partial n_{j}}%
;\quad1\leq j\leq s.
\end{equation}
The matrix elements of the Jacobian $\left(  \frac{\partial\boldsymbol{N}%
}{\partial\boldsymbol{n}}\right)  $ are
\begin{align}
\frac{\partial N_{k}}{\partial n_{j}}  &  =\frac{\partial}{\partial n_{j}%
}\left\{  n_{k}\frac{\partial\ln S_{M}}{\partial n_{k}}\right\}  =\delta
_{jk}\frac{\partial\ln S_{M}}{\partial n_{k}}+n_{k}\frac{\partial^{2}\ln
S_{M}}{\partial n_{j}\partial n_{k}}\nonumber\\
&  =\delta_{jk}\frac{S_{M-1}\left(  \jmath\right)  }{S_{M}}-n_{k}\frac
{S_{M-1}\left(  \jmath\right)  S_{M-1}\left(  k\right)  }{S_{M}^{2}}%
+n_{k}\frac{S_{M-2}\left(  jk\right)  }{S_{M}}.
\end{align}


\section{Discussion\label{Discussion}}

\appendix


\section{State Density and Reduced Density Operators \label{Denmat}}

The convex set $S_{N}$ of normalized $N$-particle density operators $D$ is
defined as
\begin{equation}
S_{N}=\left\{  D\left\vert D\in\mathcal{B}_{1+}\left(  \mathcal{H}^{N}\right)
;\operatorname{Tr}D=1\right.  \right\}  ,
\end{equation}
where $\mathcal{B}_{1+}\left(  \mathcal{H}^{N}\right)  $ is the cone of
positive trace class operators, defined as
\begin{equation}
\mathcal{B}_{1+}\left(  \mathcal{H}^{N}\right)  =\left\{  D\in\mathcal{B}%
_{1}\left(  \mathcal{H}^{N}\right)  \left\vert \overset{}{\left\langle
\Psi\left\vert D\Psi\right.  \right\rangle }\geq0\quad\forall\Psi
\in\mathcal{H}^{N}\right.  \right\}  ,
\end{equation}
where $\mathcal{B}_{1}\left(  \mathcal{H}^{N}\right)  $ is the space of trace
class operators defined as
\begin{equation}
\mathcal{B}_{1}\left(  \mathcal{H}^{N}\right)  =\left\{  \left.  X\right\vert
\operatorname{Tr}\left\{  X^{\dagger}X\right\}  <\infty\right\}
\end{equation}
which is contained in the space of bounded $N$ electron operators
$\mathcal{B}\left(  \mathcal{H}^{N}\right)  $ defined as
\begin{equation}
\mathcal{B}\left(  \mathcal{H}^{N}\right)  =\left\{  X\left\vert \sup_{\Psi
\in\mathcal{H}^{N}}\left\vert \frac{\left\langle \Psi\left\vert X\Psi\right.
\right\rangle }{\left\langle \Psi\left\vert \Psi\right.  \right\rangle
}\right\vert <\infty\right.  \right\}  .
\end{equation}
The convex set property of $S_{N}$ is explicitly displayed by
\begin{equation}
S_{N}=\left\{  \left.  \sum_{1\leq j\leq m}\alpha_{j}D_{j}\right\vert \sum
_{j}\alpha_{j}=1;0\leq\alpha_{j}\leq1,D_{j}\in S_{N},1\leq j\leq m\right\}  .
\end{equation}
A set $C$ is a cone if $X\in C\Leftrightarrow\alpha X\in C$ $\forall\alpha>0.$

The spaces $\mathcal{B}\left(  \mathcal{H}^{N}\right)  $ and $\mathcal{B}%
_{1}\left(  \mathcal{H}^{N}\right)  $ are Banach spaces with norms defined by
\begin{equation}
\left\Vert X\right\Vert =\sup_{\Psi\in\mathcal{H}^{N}}\left\vert
\frac{\left\langle \Psi\left\vert X\Psi\right.  \right\rangle }{\left\langle
\Psi\left\vert \Psi\right.  \right\rangle }\right\vert
\end{equation}
and
\begin{equation}
\left\Vert X\right\Vert _{1}=\operatorname{Tr}\left\{  \left(  X^{\dagger
}X\right)  ^{\frac{1}{2}}\right\}
\end{equation}
respectively. A closely related space that we have utilized is the space of
Hilbert Schimdt operators $\mathcal{B}_{2}\left(  \mathcal{H}^{N}\right)  $
defined as
\begin{equation}
\mathcal{B}_{2}\left(  \mathcal{H}^{N}\right)  =\left\{  X\left\vert \left\{
\operatorname{Tr}\left\{  X^{\dagger}X\right\}  \right\}  ^{\frac{1}{2}%
}<\infty\right.  \right\}  ,
\end{equation}
which is a Hilbert space with inner product given by
\begin{equation}
\left(  X\left\vert Y\right.  \right)  =\operatorname{Tr}\left\{  X^{\dagger
}Y\right\}  .
\end{equation}


\section{Second Quantized Density Operators \label{SecQuant}}

Fermion second quantized operators $\left\{  \left.  a_{j}^{\dagger
}\right\vert 1\leq j\leq r_{1}\right\}  $ are defined by the action on the
vacuum state $\left\vert \phi\right\rangle $ by
\begin{equation}
a_{j}^{\dagger}\left\vert \phi\right\rangle =\left\vert \varphi_{j}%
\right\rangle ,1\leq j\leq r_{1}%
\end{equation}
where $\left\{  \left.  \left\vert \varphi_{j}\right\rangle \right\vert 1\leq
j\leq r_{1}\right\}  $ is a complete orthonormal basis of one-electron space
$\mathcal{H}^{1}$ and the anti-commutation relationships
\begin{equation}
\left[  a_{j}^{\dagger},a_{k}\right]  _{+}=a_{j}^{\dagger}a_{k}+a_{k}%
a_{j}^{\dagger}=\delta_{jk};1\leq j,k\leq r_{1}.
\end{equation}
The field operators $\left\{  \left.  \Phi\left(  \boldsymbol{r}\right)
^{\dagger}\right\vert \boldsymbol{r\in}\mathbb{R}^{3}\right\}  $ are defined
by
\begin{equation}
\Phi\left(  \boldsymbol{r}\right)  ^{\dagger}=\sum_{1\leq j\leq r_{1}}%
\varphi_{j}\left(  \boldsymbol{r}\right)  a_{j}^{\dagger}%
\end{equation}
and the density operators by
\begin{equation}
\Phi\left(  \boldsymbol{r}\right)  ^{\dagger}\Phi\left(  \boldsymbol{r}%
\right)  =\sum_{1\leq j,k\leq r_{1}}\varphi_{j}\left(  \boldsymbol{r}\right)
\varphi_{k}\left(  \boldsymbol{r}\right)  ^{\ast}a_{j}^{\dagger}a_{k}.
\end{equation}


\section{Regular Points \label{Regular}}

The tangent space $\mathcal{T}_{A}$ at the point $A$ of the constraint
$\Xi_{2}$ is given by
\begin{equation}
\mathcal{T}_{A}=\left\{  B\left\vert \left(  A\left\vert \Phi\left(
\boldsymbol{r}\right)  ^{\dagger}\Phi\left(  \boldsymbol{r}\right)  \right.
B\right)  =0\forall\boldsymbol{r}\right.  \right\}  \subseteq\ker d\Xi
_{2}\left(  A\right)
\end{equation}
and the normal space $\mathcal{N}_{A}$ by
\begin{equation}
\mathcal{N}_{A}=\left\{  \left.  B\right\vert \left(  A\left\vert \Phi\left(
\boldsymbol{r}\right)  ^{\dagger}\Phi\left(  \boldsymbol{r}\right)  \right.
B\right)  \neq0\forall\boldsymbol{r\in}\mathbb{R}^{3}\right\}
\end{equation}
One can show that $\left\{  \left.  \left[  d\Xi_{2}\left[  \boldsymbol{r}%
\right]  \left(  A\right)  \right]  \left(  \cdot\right)  \equiv\left(
\Phi\left(  \boldsymbol{r}\right)  ^{\dagger}\Phi\left(  \boldsymbol{r}%
\right)  A\right\vert \equiv\operatorname{Tr}\left\{  A\Phi\left(
\boldsymbol{r}\right)  ^{\dagger}\Phi\left(  \boldsymbol{r}\right)
\cdot\right\}  \right\vert \boldsymbol{r\in}\mathbb{R}^{3}\right\}  $ is a
generalized basis, (see Appendix basis \ref{Bases}), for $\mathcal{N}_{A}$
which shows that
\begin{equation}
\mathcal{B}_{2sa}\left(  \mathcal{H}^{N}\right)  =\mathcal{T}_{A}%
\oplus\mathcal{N}_{A}%
\end{equation}
and
\begin{equation}
\left[  d\Xi_{2}\left(  A\right)  \right]  :\mathcal{B}_{2sa}\left(
\mathcal{H}^{N}\right)  \rightarrow L_{1}\left(  \mathbb{R}^{3}\right)
\end{equation}
is surjective so that $A$ is a \emph{regular }\cite{Hestenes}\emph{ }point of
the constraint.

\section{Bases \label{Bases}}

\section{Super Operators \label{super}}

Super operators, which could be unbounded, and are defined by the quadratic
form
\begin{equation}
\left(  X\left\vert \widehat{T}Y\right.  \right)  =\operatorname{Tr}\left\{
XTY\right\}  \,\forall X\in\mathcal{B}_{2sa}\left(  \mathcal{H}^{N}\right)
,\,Y\in\operatorname{Dom}\left\{  \widehat{T}\right\}  ,
\end{equation}
where
\[
\operatorname{Dom}\left\{  \widehat{T}\right\}  =\left\{  \left.  Y\right\vert
\operatorname{Tr}\left\{  XTY\right\}  <\infty\,\forall X\in\mathcal{B}%
_{2sa}\left(  \mathcal{H}^{N}\right)  \right\}  .
\]


\section{State Normalization \label{Normalization}}

If $D^{N}\in\mathcal{B}_{2}\left(  \mathcal{H}^{N}\right)  $ then
\begin{equation}
\operatorname{Tr}\left\{  \mathcal{N}D^{N}\right\}  =N\operatorname{Tr}%
\left\{  D^{N}\right\}  ,
\end{equation}
where
\begin{equation}
\mathcal{N}=\int\Phi\left(  \boldsymbol{r}\right)  ^{\dagger}\Phi\left(
\boldsymbol{r}\right)  d\boldsymbol{r,}%
\end{equation}
as
\begin{equation}
\mathcal{N}\left\vert \Psi\right\rangle =N\left\vert \Psi\right\rangle
\quad\forall\left\vert \Psi\right\rangle \in\mathcal{H}^{N}.
\end{equation}
For any state $D$
\begin{equation}
\operatorname{Tr}\left\{  \mathcal{N}D\right\}  =\int\rho\left(
\boldsymbol{r,\xi}\right)  d\boldsymbol{r}d\boldsymbol{\xi,}%
\end{equation}
thus
\begin{equation}
N\operatorname{Tr}\left\{  D^{N}\right\}  =N\Rightarrow\operatorname{Tr}%
\left\{  D^{N}\right\}  =1.
\end{equation}


\section{Complex Conjugate Variables \label{ComplexConjugate}}

\section{Spin\label{Spin}}

\section{Symmetric Function Relationships}

\bibliographystyle{apsrev}
\bibliography{agpopt,bw,novel,San2008,coleman,Coleman2,bibrefs,lowdin-per-olov,Janak}

\end{document}